\documentclass[11pt]{amsart}

\usepackage{amsmath}
\usepackage{amssymb}
\usepackage{amsthm}
\usepackage{booktabs}
\usepackage{array}
\usepackage{hyperref}

% ---------------------------------------------------------------------------
% Theorem environments
% ---------------------------------------------------------------------------
\newtheorem{theorem}{Theorem}[section]
\newtheorem{proposition}[theorem]{Proposition}
\newtheorem{lemma}[theorem]{Lemma}
\newtheorem{corollary}[theorem]{Corollary}
\theoremstyle{definition}
\newtheorem{definition}[theorem]{Definition}
\theoremstyle{remark}
\newtheorem{remark}[theorem]{Remark}

\numberwithin{equation}{section}
\setlength{\emergencystretch}{2em}

% ---------------------------------------------------------------------------
% Macros
% ---------------------------------------------------------------------------
\newcommand{\NN}{\mathbb{N}}
\newcommand{\ZZ}{\mathbb{Z}}
\newcommand{\QQ}{\mathbb{Q}}
\newcommand{\RR}{\mathbb{R}}
\newcommand{\KK}{\mathbb{K}}
\newcommand{\cK}{\mathcal{K}}
\newcommand{\cF}{\mathcal{F}}
\newcommand{\cM}{\mathcal{M}}
\newcommand{\cP}{\mathcal{P}}
\newcommand{\fp}{\mathfrak{p}}
\newcommand{\fl}{\mathfrak{l}}
\newcommand{\fP}{\mathfrak{P}}
\newcommand{\fL}{\mathfrak{L}}
\newcommand{\one}{\mathbf{1}}
\newcommand{\eps}{\varepsilon}
\newcommand{\ind}[1]{[\,#1\,]}
\newcommand{\lev}{\operatorname{lev}^{*}}
\newcommand{\adj}{\operatorname{adj}}
\newcommand{\Li}{\operatorname{Li}}
\newcommand{\cross}{\operatorname{cross}}
\newcommand{\Vand}{\operatorname{Vand}}
\newcommand{\artanh}{\operatorname{artanh}}
\newcommand{\lean}[1]{{\def\_{\textunderscore\allowbreak}\texttt{#1}}}

\begin{document}

\title{\texorpdfstring{$\pi^2-18\log^2\varphi$ is irrational}{pi^2 - 18 log^2 phi is irrational}}

\author{Shawn Krisman}

\subjclass[2020]{Primary 11J72; Secondary 33B30, 11R27, 31A15, 41A21, 68V20}
\keywords{golden ratio, Legendre chi function, dilogarithm, irrationality, Hankel determinants, units of quadratic
fields, logarithmic potential theory, formal verification, Lean}

\date{\today}

\begin{abstract}
Let $\varphi=(1+\sqrt5)/2$. We prove that $\pi^2-18\log^2\varphi$ is irrational. Equivalently, the value
$\chi_2(\sqrt5-2)=\pi^2/24-\frac34\log^2\varphi$ of Legendre's chi function is irrational. The proof transposes our
proof that $\chi_2(1/3)$ is irrational to the point $z=\sqrt5-2$, a unit of norm $-1$ in $\QQ(\sqrt5)$. We take
poles at the even integers $\pm2j$, $|j|\le h/2$, and atoms at the odd positive integers. A reflection $t\mapsto-t$,
combined with the Galois conjugation $z\mapsto-1/z$, shows that the Hankel determinants are values of polynomials
$\Delta_h(X)\in\QQ[X]$, although their entries lie in $\QQ(\sqrt5)$; so the conjugate place, where the defining series
diverges, costs nothing. Because $z$ is a unit, there are no special primes. The odd primes cost at most $4/3$ per
$h^2$, through a bordered Lagrange Gram identity and a class count. The prime $2$ costs nothing at the scale $h^2$, because
$1-z^2=4z$. The size is bounded by a discrete log-gas argument against a dual certificate checked in interval
arithmetic, which gives at most $-1.3592$ per $h^2$; nonvanishing comes from a reduction modulo a prime $\ell=(h+1)/3$. The
margin is $-0.02588\ldots$ per $h^2$. The theorem is formalised in Lean~4 with Mathlib, using only Lean's three standard axioms; the
formal proof of the analytic bound follows a different route from the one written here. We also prove a two-point
statement at units of a family of biquadratic fields: for $n\in\{730,740,750,760,800,1000\}$, either
$\chi_2(z)+\chi_2(w)\notin\QQ$ or $\chi_2(z)-\chi_2(w)\notin\sqrt{n^2-16}\,\QQ$, where $z>w$ are the two roots in $(0,1)$ of
$X^4+nX^3+2X^2-nX+1$. This second result is computer-assisted and not formalised.
\end{abstract}

\maketitle

\tableofcontents

% ===========================================================================
\section{Introduction}\label{sec:intro}
% ===========================================================================

Let $\varphi=(1+\sqrt5)/2$ be the golden ratio, and for $|q|\le1$ let
\[
  \chi_2(q)=\sum_{n=0}^{\infty}\frac{q^{2n+1}}{(2n+1)^2}=\tfrac12\bigl(\Li_2(q)-\Li_2(-q)\bigr)
\]
be Legendre's chi function \cite[\S1.8]{Lewin1981}.

\begin{theorem}[Theorem U]\label{thm:main}
The number
\[
  \pi^2-18\log^2\varphi=5.70143763069985955388\ldots
\]
is irrational. Equivalently, the number
\[
  \chi_2(\sqrt5-2)=\frac{\pi^2}{24}-\frac34\log^2\varphi=0.23755990127916081474\ldots
\]
is irrational.
\end{theorem}

The two forms are equivalent by the identity
\begin{equation}\label{eq:closedform}
  \chi_2(\sqrt5-2)=\frac{\pi^2}{24}-\frac34\log^2\varphi ,
\end{equation}
which follows from Landen's identity
$\chi_2\bigl(\tfrac{1-x}{1+x}\bigr)+\chi_2(x)=\tfrac{\pi^2}{8}+\tfrac12\log x\,\log\tfrac{1+x}{1-x}$
\cite[App.~A.2.1, (24)]{Lewin1981} at $x=\sqrt5-2=\varphi^{-3}$, for which $(1-x)/(1+x)=\varphi^{-1}$, together with the
classical value $\chi_2(\varphi^{-1})=\pi^2/12-\frac34\log^2\varphi$. The formalisation proves \eqref{eq:closedform}
from these identities (Section~\ref{sec:lean}).

\subsection{Context}\label{sec:context}
% Related work from schanuel/four_exp/strong4/pilot_U/novelty/NOVELTY.md (finished version, Sections 2--5, read
% 2026-11-07 by the whole-chain referee; its Section 3 "open search" is complete and adds no relevant prior result).
We have not found Theorem~\ref{thm:main}, or the irrationality of any number $\pi^2+c\log^2\varphi$ with
$c\in\QQ\setminus\{0\}$, in the literature, and as far as we can see it does not follow from known results. By the
analytic class number formula, $L(1,\chi_5)=2\log\varphi/\sqrt5$ for the quadratic character $\chi_5$ modulo $5$, so the
theorem can also be written as $6\zeta(2)-\frac{45}{2}L(1,\chi_5)^2\notin\QQ$.
\begin{itemize}
\item \emph{The golden-ratio dilogarithm values are different numbers.} The classical evaluations
  \[
    \begin{aligned}
    \Li_2(\varphi^{-2})&=\tfrac{\pi^2}{15}-\log^2\varphi,&\quad \Li_2(\varphi^{-1})&=\tfrac{\pi^2}{10}-\log^2\varphi,\\
    \Li_2(-\varphi^{-1})&=-\tfrac{\pi^2}{15}+\tfrac12\log^2\varphi,&\quad \Li_2(-\varphi)&=-\tfrac{\pi^2}{10}-\log^2\varphi
    \end{aligned}
  \]
  (Landen; see \cite{Lewin1981,Lewin1982}) all lie in $\QQ\pi^2+\QQ\log^2\varphi$. Normalised to the form
  $\pi^2+c\log^2\varphi$ they have $c=-15$, $-10$, $-\frac{15}{2}$ and $+10$, while Theorem~\ref{thm:main} has $c=-18$.
  Since $\log\varphi$ is transcendental, $\pi^2+c\log^2\varphi$ can be algebraic for at most one $c\in\overline\QQ$; this
  does not say which. The irrationality of one member of the family says nothing about another, and as far as we know
  the irrationality of each of the four classical values is open.
\item \emph{Homogeneous statements.} The theorem of Gel'fond and Schneider \cite[Ch.~2]{Baker1975}, applied to
  $e^{\log\varphi}=\varphi$ and $e^{i\pi}=-1$, shows that $\log\varphi/(i\pi)$ is transcendental, so
  $\log^2\varphi/\pi^2\notin\overline{\QQ}$. This says that $\pi^2$ and $\log^2\varphi$ are linearly independent over
  $\QQ$; it says nothing about a relation that also involves $1$. The only member of the family
  $\pi^2+c\log^2\varphi$, $c\in\QQ$, whose irrationality we know to be established is $c=0$. The linear independence of
  $1,\pi^2,\log^2\varphi$ over $\QQ$ is open as far as we know.
\item \emph{Criteria at algebraic points.} The linear independence criteria for polylogarithms and Lerch functions at
  algebraic points \cite{DHK2020,DHK2022,DHK2022b,Kawashima2026}, and the explicit form of the Bombieri--Andr\'e
  principle of global relations \cite{Orr2026}, need a point that is small in comparison with its height. Here
  $\chi_2(\sqrt5-2)=\frac12(\Li_2(z)-\Li_2(-z))=z+\frac z4\Phi_2(\frac12,z^2)$, where
  $\Phi_2(x,q)=\sum_{k\ge0}q^{k+1}/(k+x+1)^2$ is the Lerch function and $z=\sqrt5-2=\varphi^{-3}$ is a unit of height
  $\frac32\log\varphi=0.7218\ldots$ whose conjugate $-\varphi^3$ has absolute value $>1$, so that $|z|\,|\sigma z|=1$.
  Evaluating the criteria in the form they are stated, their key quantity is negative by a wide margin at this
  point: about $-26.1$ for \cite[Thm.~2.2]{Kawashima2026}, and about $-15.1$ and $-14.5$ for the two ways of applying
  \cite[Thm.~1]{DHK2022b} (through $\Li_2(\pm z)$, or through $\Phi_2(\frac12,z^2)$). The
  Bombieri--Andr\'e bound of \cite{Orr2026}, applied to this system, needs by our estimate a height of about $1450$.
  These evaluations are ours and have not been checked by the authors of the criteria.
\item \emph{Rational points.} The irrationality and linear independence results for dilogarithms at rational points
  \cite{Hata1990,Hata1993,RhinViola2005,RhinViola2019,ViolaZudilin2018}, the recent claims
  \cite{Prellberg2026,TanZhang2026} (unrefereed), and our result $\chi_2(1/3),\chi_2(1/5)\notin\QQ$ \cite{Chi2Paper} all
  concern rational arguments, as do the weight-two results of Calegari, Dimitrov and Tang \cite{CDT} (at sixth roots
  of unity, and for products $\log(1+1/m)\log(1+1/n)$) and Hata's irrationality of $\log(1+1/q)\log(1-1/q)$
  \cite{Hata1998}. The results on products of logarithms are homogeneous and contain no $\pi^2$.
\item \emph{The closest analogue.} The nearest result we know is $\Li_2(\frac12)=\frac{\pi^2}{12}-\frac12\log^22\notin\QQ$,
  proved in Lean by Q.~Tang \cite{TangLi2} and, with many more rational points, by S.~Liu \cite{LiuLi2}; neither is
  accompanied by a refereed paper. It has the same shape, $\pi^2+c\log^2$ of a number, with $c=-6$ and the number $2$
  instead of a unit. Over $\QQ(\sqrt5)$ we have found nothing comparable.
\end{itemize}
Our searches for prior work, in four passes, covered the web, arXiv (through its API), zbMATH Open, OpenAlex, Crossref,
Semantic Scholar (only partly, because of rate limits), the Hexagon repository, the \texttt{openai/math} manuscript list
and recent Lean repositories (the last pass on 7 November 2026). They did not cover MathSciNet or Google Scholar, and no
expert has been consulted. The searches were carried out by AI systems.

\subsection{The method}\label{sec:method}
Put
\[
  z:=\sqrt5-2=\varphi^{-3},\qquad y:=z^2,\qquad T:=\chi_2(z)=\sum_{t\text{ odd}\ge1}\frac{z^t}{t^2}.
\]
Then $z^2+4z-1=0$, so $1-y=4z$, and the Galois conjugate of $z$ is $-1/z$: $z$ is a unit of norm $-1$ in
$\KK:=\QQ(\sqrt5)$. The proof transposes the proof of \cite{Chi2Paper} for $\chi_2(1/q)$ to this point. For even $h$ we
take the $K:=h+1$ poles $x_m:=2(m-h/2)$, i.e.\ the even integers in $[-h,h]$, put $D(t):=\prod_m(t-x_m)$, and consider
the functional
\[
  \Lambda(g):=\sum_{t\text{ odd}\ge1}z^t\Bigl(\frac{g}{D}\Bigr)'(t)\qquad(g\in\QQ[t]).
\]
Partial fractions give $\Lambda=A+T\,B$ with $\KK$-valued functionals $A,B$, so the Hankel determinant
$\det(\Lambda(t^{i+j}))_{i,j<h}$ is the value at $X=T$ of $\Delta_h(X):=\det(A(t^{i+j})+XB(t^{i+j}))_{i,j<h}$, a
polynomial of degree at most $h$.

What differs from \cite{Chi2Paper} is the \emph{reflection} (Lemma~\ref{lem:R}). The pole set is symmetric under $t\mapsto-t$, $D(-t)=-D(t)$,
and the conjugation $\sigma\colon z\mapsto-1/z$ maps the partial sums of $\chi_2$ at the poles $x$ to those at $-x$.
Together these give $\sigma(\Lambda_X(t^m))=(-1)^m\Lambda_X(t^m)$, which forces $\Delta_h\in\QQ[X]$ although its entries
lie in $\KK$. So we never pay for the conjugate place, where $|\sigma z|=\varphi^3>1$ and the defining series diverges.
The three inputs are then as in \cite{Chi2Paper}.
\begin{enumerate}
\item \emph{Denominators} (Sections~\ref{sec:odd}--\ref{sec:two}). Because $z$ is a unit, the moments of $\Lambda$ only
  have the denominators $(1-y)^{b+1}=(4z)^{b+1}$, so there are no special primes, and no gain prime either. Every odd
  prime, including $5$, satisfies the same bounds. Primes $p<K/3$ obey a plateau bound $2h\lev(p)+v_p(h!)$, which we
  prove through a bordered version of the Lagrange Gram identity (Lemmas~\ref{lem:bordered} and~\ref{lem:nine}), and
  primes $K/3\le p<K$ obey a class count $3(K-p)$ (Lemmas~\ref{lem:eight} and~\ref{lem:classcount}). The prime number
  theorem sums these to $\frac23+\frac23=\frac43$ per $h^2$ (Theorem~\ref{thm:A}). The prime $2$ costs only
  $O(h\log h)$ (Lemma~\ref{lem:ten}), because $v_2(1-y)=2$.
\item \emph{Size} (Section~\ref{sec:analytic}). A discrete log-gas argument bounds $h^{-2}\log|\Delta_h(T)|$ by the value
  $\Phi$ of any dual certificate for a constrained equilibrium problem with external field
  $x\log z-\frac12\int_{x-1}^{x+1}\log|s|\,ds$ (Theorem~\ref{thm:B}). An explicit certificate, checked in interval
  arithmetic, gives $\Phi\le-1.3592161$ (Theorem~\ref{thm:cert}).
\item \emph{Nonvanishing} (Section~\ref{sec:nonvan}). For $h=3\ell-1$ with $\ell\ge7$ prime, the Gram matrix of
  $\Lambda$ in the Lagrange basis is block diagonal modulo a prime of $\KK$ above $\ell$, with $3\times3$ blocks in
  closed form, so $\Delta_h(a/b)\ne0$ whenever $\ell^2\nmid b$ (Lemma~\ref{lem:P}). Unlike in \cite{Chi2Paper}, every
  prime $\ell\ge7$ works.
\end{enumerate}
Table~\ref{tab:budget} collects the contributions.

\begin{table}[ht]
\centering
\begin{tabular}{@{}llr@{}}
\toprule
contribution (per $h^2$, $\limsup$) & source & value\\
\midrule
odd primes, including $5$ & Thm.~\ref{thm:A} & $4/3=1.3333333$\\
$p=2$ & Lemma~\ref{lem:ten} & $0$\\
special primes & ($z$ is a unit) & none\\
\midrule
arithmetic & & $1.3333333$\\
analytic $\Phi$ & Thms.~\ref{thm:B}, \ref{thm:cert} & $\le-1.3592161$\\
\midrule
margin & & $\le-0.025882$\\
\bottomrule
\end{tabular}
\caption{The exponential budget. The analytic value is the upper endpoint $-1.3592161834586374$ of the certified Arb
ball, rounded up; the margin is $-0.0258828501\ldots$, rounded up.}\label{tab:budget}
\end{table}

The margin is about $1.9\%$ of the arithmetic part. It is smaller than in \cite{Chi2Paper} for $q=3$ in relative terms,
but no constant in the budget is estimated: each is a proved upper bound.

\subsection{The formalisation}
Theorem~\ref{thm:main} and the identity \eqref{eq:closedform} are formalised in Lean~4 \cite{LeanFour} with Mathlib
\cite{Mathlib} and checked by the Lean kernel. The headline theorems are \lean{Chi2U.chi2\_sqrt5\_sub\_two\_irrational}
($\chi_2(\sqrt5-2)\notin\QQ$), \lean{Chi2U.golden\_irrational'} ($\pi^2-18\log^2\varphi\notin\QQ$) and
\lean{Chi2U.chi2\_sqrt5\_sub\_two\_eq} (the identity \eqref{eq:closedform}). They depend only on Lean's three standard
axioms (\lean{propext}, \lean{Classical.choice}, \lean{Quot.sound}); there is no \lean{sorry} and no
\lean{native\_decide}. The formal proof of the analytic bound follows the lattice-energy route of \cite{Chi2Paper}, not
the route of Section~\ref{sec:analytic}, and uses its own, coarser certificate; Section~\ref{sec:lean} describes both
differences.

\subsection{Two points}
Section~\ref{sec:M1} transposes the proof once more, to units of the biquadratic fields $\QQ(z)$ where $z$ is a root of
$X^4+nX^3+2X^2-nX+1$. There the reflection removes the two conjugates $-1/z$ and $-1/w$, where the series diverge,
and the remaining conjugate $w$ is a second convergent place, which enters through a norm. The result is a disjunction, proved separately for six values of $n$,
with interval-arithmetic certificates. The smallest instance is the following.

\begin{theorem}[M1 at $n=730$]\label{thm:M1-730}
Let $z>w$ be the two roots in $(0,1)$ of $X^4+730X^3+2X^2-730X+1$. Then at least one of the following holds:
\[
  \frac{\pi^2}{8}-\frac18\log^2\bigl(366+\sqrt{133955}\bigr)+\frac18\log^2\bigl(364+11\sqrt{1095}\bigr)\notin\QQ,
\]
\[
  \chi_2(z)-\chi_2(w)\notin\sqrt{1101}\,\QQ .
\]
\end{theorem}

The first number is $\chi_2(z)+\chi_2(w)$, by Landen's identity. Theorem~\ref{thm:M1-730} is not formalised.

\subsection{Organisation and notation}
Section~\ref{sec:construction} sets up the functional, proves the pole expansion and the reflection, and states the
criterion. Sections~\ref{sec:odd} and~\ref{sec:two} bound the denominators, Section~\ref{sec:analytic} bounds the size,
Section~\ref{sec:nonvan} proves nonvanishing, and Section~\ref{sec:proof} assembles the proof. Section~\ref{sec:M1}
treats two points, Section~\ref{sec:lean} describes the formalisation, and Section~\ref{sec:comp} lists the computations.

Throughout, $p$ and $\ell$ denote primes, $v_p$ is the $p$-adic valuation, $\lfloor\cdot\rfloor$ is the floor,
$(x)_+=\max(x,0)$, and $\ind{P}$ is $1$ if $P$ holds and $0$ otherwise. For a function $G$ on $\ZZ_{\ge0}$,
$\Delta G(s)=G(s+1)-G(s)$ is the forward difference. We write $\theta(x)=\sum_{p\le x}\log p$ and
$\psi(x)=\sum_{p^e\le x}\log p$, and $s_2(b)$ for the binary digit sum of $b$, so $v_2(b!)=b-s_2(b)$.

% ===========================================================================
\section{The construction}\label{sec:construction}
% ===========================================================================

\subsection{Data}\label{sec:data}
Let $\KK=\QQ(\sqrt5)$, embedded in $\RR$ with $\sqrt5>0$, and let $\sigma$ be its nontrivial automorphism. With
$z=\sqrt5-2$ and $y=z^2$ we have
\[
  z^2+4z-1=0,\qquad 1-y=4z,\qquad \sigma(z)=-1/z .
\]
So $z$, $y$ and $1-y$ are units at every odd prime, including $5$, and $1-y$ has $2$-adic valuation $2$.

Fix an even $h\ge2$ and put $J:=h/2$ and $K:=h+1$. The poles are $x_m:=2(m-J)$ for $0\le m<K$, i.e.\ the even integers
$\pm2j$, $|j|\le J$, and
\begin{gather*}
  D(t):=\prod_{m<K}(t-x_m),\\
  D'(x_m)=\prod_{l\ne m}2(m-l)=(-1)^{K-1-m}2^{K-1}m!\,(K-1-m)! .
\end{gather*}
For a rational function $f$ regular at the odd positive integers and of polynomial growth put
\[
  \cM(f):=\sum_{t\text{ odd}\ge1}z^tf'(t),
\]
an absolutely convergent series, and $\Lambda(g):=\cM(g/D)$ for $g\in\QQ[t]$. For an even integer $x$ define the partial
sums
\[
  P(x):=\begin{cases}\displaystyle\sum_{n\text{ odd},\,0<n<|x|}\frac{z^n}{n^2},&x<0,\\[10pt]
  \displaystyle-\sum_{n\text{ odd},\,-x<n<0}\frac{z^n}{n^2},&x>0,\\[8pt] 0,&x=0.\end{cases}
\]
Only odd $n$ with $|n|\le|x|-1\le h-1$ occur.

\subsection{The pole expansion}

\begin{lemma}[Lemma~1$'$, pole expansion]\label{lem:pole}
For every even integer $x$, $\cM\bigl(1/(t-x)\bigr)=-z^x\bigl(T-P(x)\bigr)$. For a polynomial $F$,
\begin{equation}\label{eq:mahler}
  \cM(F)=z\sum_{b\ge0}\Delta^bG(0)\,\frac{y^b}{(1-y)^{b+1}},\qquad G(s):=F'(2s+1),
\end{equation}
a finite sum. Consequently, writing $g/D=R_g+\sum_mg(x_m)/(D'(x_m)(t-x_m))$ with $R_g\in\QQ[t]$,
\begin{gather*}
  \Lambda(g)=A(g)+T\,B(g),\\
  A(g):=\cM(R_g)+\sum_m\frac{g(x_m)}{D'(x_m)}z^{x_m}P(x_m),\qquad
  B(g):=-\sum_m\frac{g(x_m)}{D'(x_m)}z^{x_m},
\end{gather*}
with $A,B\colon\QQ[t]\to\KK$ linear.
\end{lemma}

\begin{proof}
$\cM(1/(t-x))=-\sum_{t\text{ odd}\ge1}z^t/(t-x)^2$. With $n=t-x$ (odd, $n\ge1-x$) and $z^t=z^xz^n$, the sum is
$z^x\sum_{n\text{ odd}\ge1-x}z^n/n^2=z^x(T-P(x))$ by the definition of $P$. For a polynomial, substitute $t=2s+1$,
expand $G(s)=\sum_b\Delta^bG(0)\binom sb$, and use $\sum_{s\ge0}\binom sby^s=y^b/(1-y)^{b+1}$. All series converge
absolutely since $0<z<1$.
\end{proof}

For an indeterminate $X$ put $\Lambda_X:=A+X\,B$, $E_m:=\Lambda_X(t^m)\in\KK[X]$ and
\begin{gather*}
  \Delta_h(X):=\det\bigl(E_{i+j}\bigr)_{0\le i,j<h}\in\KK[X],\qquad\deg\Delta_h\le h,\\
  \Delta_h(T)=\det\bigl(\Lambda(t^{i+j})\bigr)_{i,j<h}.
\end{gather*}
We use the same notation for every finite family: if $a_0,\dots,a_{h-1}$ and $b_0,\dots,b_{h-1}$ have degrees $<h$ and
coefficient matrices $C_a,C_b$, then $\det(\Lambda_X(a_ib_j))=\det C_a\det C_b\,\Delta_h(X)$.

\subsection{The reflection}
Extend $A$ and $B$ to rational functions $f=R+\sum_mc_m/(t-x_m)$ ($R\in\QQ[t]$, $c_m\in\QQ$) by
$A(f):=\cM(R)+\sum_mc_mz^{x_m}P(x_m)$ and $B(f):=-\sum_mc_mz^{x_m}$, so that $\cM(f)=A(f)+TB(f)$, and put
$\tilde f(t):=f(-t)$.

\begin{lemma}[Lemma~R, reflection]\label{lem:R}
For every such $f$, $\sigma(A(f))+\sigma(B(f))X=-\bigl(A(\tilde f)+B(\tilde f)X\bigr)$.
\end{lemma}

\begin{proof}
We have $\tilde f=\tilde R+\sum_m(-c_m)/(t-(-x_m))$, and the pole set is symmetric. \emph{The $X$-part.} Since $x_m$ is
even, $\sigma(z^{x_m})=(-1/z)^{x_m}=z^{-x_m}$, so $\sigma(B(f))=-\sum_mc_mz^{-x_m}=-B(\tilde f)$. \emph{The pole part
of $A$.} For $x<0$, $\sigma(P(x))=\sum_{0<n<|x|}(-1)^nz^{-n}/n^2=-\sum_{-|x|<n'<0}z^{n'}/n'^2=P(-x)$ (odd $n$, $n'=-n$),
and likewise for $x>0$. Hence $\sigma\bigl(\sum_mc_mz^{x_m}P(x_m)\bigr)=\sum_mc_mz^{-x_m}P(-x_m)$, which is minus the pole
part of $A(\tilde f)$. \emph{The polynomial part.} For a polynomial $R$, $\cM(R)=M_R(z)$ with
$M_R(q):=\sum_{t\text{ odd}\ge1}R'(t)q^t$, a rational function of $q$ with rational coefficients and poles only at
$q=\pm1$; so $\sigma(\cM(R))=M_R(-1/z)$. Since $R'$ is a polynomial, the reciprocity of rational generating functions
gives $\sum_{t\text{ odd}\ge1}R'(t)q^t=-\sum_{t\text{ odd}\le-1}R'(t)q^t$ as rational functions, the right side being
expanded at $q=\infty$: for $R'(t)=t^k$ apply $(q\,d/dq)^k$ to $\sum_{t\text{ odd}\ge1}q^t=q/(1-q^2)=-\sum_{t\text{ odd}\le-1}q^t$. At $q=-1/z$ the
right side is $-\sum_{u\text{ odd}\ge1}R'(-u)(-z)^{u}=\sum_{u\text{ odd}\ge1}R'(-u)z^u=-M_{\tilde R}(z)$, because
$\tilde R'(u)=-R'(-u)$. So $\sigma(\cM(R))=-\cM(\tilde R)$.
\end{proof}

\begin{corollary}[Corollary~R]\label{cor:R}
$\sigma(E_m)=(-1)^mE_m$ coefficientwise, and $\Delta_h(X)\in\QQ[X]$.
\end{corollary}

\begin{proof}
For $f=t^m/D$ we have $D(-t)=(-1)^KD(t)=-D(t)$ ($K$ is odd), so $\tilde f=(-1)^{m+1}f$, and Lemma~\ref{lem:R} gives
$\sigma(E_m)=(-1)^mE_m$. Thus $E_m\in\QQ[X]$ for even $m$ and $E_m\in\sqrt5\,\QQ[X]$ for odd $m$. The matrix
$\tilde E_{ij}:=E_{i+j}\sqrt5^{\,\ind{i\text{ odd}}-\ind{j\text{ odd}}}$ has entries in $\QQ[X]$, and
$\tilde E=\operatorname{diag}(\sqrt5^{\,\ind{i\text{ odd}}})\,(E_{i+j})\,\operatorname{diag}(\sqrt5^{-\ind{j\text{ odd}}})$,
so $\det\tilde E=\Delta_h$.
\end{proof}

\subsection{Valuations and the criterion}
For a prime $p$ and a prime $\fp\mid p$ of $\KK$ let $v_\fp$ be normalised by $v_\fp(p)=1$ (so $v(\sqrt5)=\frac12$ at
$p=5$), extended to $\KK[X]$ as the Gauss valuation (the minimum over the coefficients). For $p\ne2,5$ we write $v_p$;
since $2$ is inert and $5$ ramified in $\KK$ there is one prime above each $p\in\{2,5\}$, and on $\QQ[X]$ every $v_\fp$
restricts to $v_p$. For $\Delta_h\not\equiv0$ put
\begin{gather*}
  n_p:=-v_p(\Delta_h),\qquad F_h:=\Delta_h\prod_pp^{n_p}\in\ZZ[X]\ \text{(primitive)},\\
  \log|F_h(T)|=\log|\Delta_h(T)|+\sum_pn_p\log p .
\end{gather*}
Upper bounds for $n_p$ may be computed from the $\KK$-valued entries at $\fp\mid p$, since $\Delta_h\in\QQ[X]$.

\begin{proposition}[Lemma~E]\label{prop:criterion}
Let $\xi\in\RR$ and let $\mathcal H$ be an infinite set of positive integers. Suppose that for $h\in\mathcal H$ there are
$F_h\in\ZZ[X]$ with $\deg F_h\le h$ and $\limsup_{h\in\mathcal H}h^{-2}\log|F_h(\xi)|<0$, and that $F_h(a/b)\ne0$ for
all $h\in\mathcal H$, where $a\in\ZZ$, $b\ge1$. Then $\xi\ne a/b$.
\end{proposition}

\begin{proof}
If $\xi=a/b$, then $b^hF_h(a/b)$ is a nonzero integer, so $h\log b+\log|F_h(\xi)|\ge0$ for $h\in\mathcal H$. Dividing by
$h^2$ contradicts the hypothesis.
\end{proof}

This is the criterion of \cite[Prop.~3.1]{Chi2Paper}, with the nonvanishing set allowed to depend on $b$.

% ===========================================================================
\section{Denominators I: the odd primes}\label{sec:odd}
% ===========================================================================

Throughout this section $p$ is an odd prime and $v=v_\fp$ for a prime $\fp\mid p$ of $\KK$. Since $D$ is monic with
integer coefficients, $R_g\in\ZZ[t]$ for $g\in\ZZ[t]$, and by \eqref{eq:mahler}, with $y$ and $1-y=4z$ units at $\fp$,
\begin{equation}\label{eq:step3}
  v\bigl(\cM(F)\bigr)\ge\min_{t\text{ odd}\ge1}v_p\bigl(F'(t)\bigr)\qquad(F\in\ZZ_{(p)}[t]),
\end{equation}
because $\Delta^bG(0)$ is an integer combination of values $F'(t)$ at odd $t$. For odd $t$ we have
$t-x_m=(t+2J)-2m$, so for a level $p^e$ the indices $m$ with $p^e\mid t-x_m$ form one residue class $m_0(t)$ modulo
$p^e$ among the $K$ consecutive integers $0,\dots,K-1$. Put
\[
  \lev(p):=\#\{e\ge1:p^e\le h\}.
\]

\subsection{The atom expansion and the class count}
Let $V_{m,i}:=x_m^i$ ($m<K$, $i<h$), $\cP_{ij}:=\cM(R_{t^{i+j}})$ and $\gamma_m(X):=-z^{x_m}(X-P(x_m))/D'(x_m)$. By
Lemma~\ref{lem:pole}, $(E_{i+j})=\cP+V^{\mathsf T}\operatorname{diag}(\gamma)V$.

\begin{lemma}[Lemma~7$'$, atom expansion]\label{lem:atom}
\[
  \Delta_h(X)=\sum_{|I|=|I'|=|S|}\pm\det\cP_{I^cI'^c}\,\det V_{S,I}\,\det V_{S,I'}\prod_{m\in S}\gamma_m(X),
\]
over $I,I'\subseteq\{0,\dots,h-1\}$ and $S\subseteq\{0,\dots,K-1\}$. Moreover $\det V_{S,I}\in\Vand(x_S)\ZZ$, and
$\prod_{m\in S}|D'(x_m)|=\Vand(x_S)^2\cross(S)$ with $\cross(S):=\prod_{m\in S,\,l\notin S}|x_m-x_l|$.
\end{lemma}

\begin{proof}
Expand $\det(\cP+M)$ by the complementary minors of $\cP$ and $M=V^{\mathsf T}\operatorname{diag}(\gamma)V$, and each minor
$\det M_{II'}$ by the Cauchy--Binet formula. A generalised Vandermonde determinant is $\Vand(x_S)$ times a Schur
polynomial with integer coefficients. The last identity splits $\prod_{l\ne m}|x_m-x_l|$ into $l\in S$ and $l\notin S$.
\end{proof}

For an odd level $d$ put $w_d(m):=2\ind{|x_m|>d}$, let $C_c:=\{m<K:m\equiv c\bmod d\}$, $n_c:=|C_c|$, and
\[
  A_K(d):=\sum_{c\bmod d}\ \max_{S_c\subseteq C_c}\Bigl[|S_c|(n_c-|S_c|)+\sum_{m\in S_c}w_d(m)\Bigr].
\]

\begin{lemma}[Lemma~8$'$, per-level bound]\label{lem:eight}
For every odd prime $p$, $n_p\le\sum_{e\ge1,\,p^e<K}A_K(p^e)$.
\end{lemma}

\begin{proof}
$\cP$ is $\fp$-integral by \eqref{eq:step3}. In the term of Lemma~\ref{lem:atom} indexed by $(I,I',S)$, the factors
$\det V_{S,I}\det V_{S,I'}$ cancel $\Vand(x_S)^2$, so by Lemma~\ref{lem:atom} the term has valuation at least
$-v_p(\cross(S))+\sum_{m\in S}v(X-P(x_m))$, using that $z^{x_m}$ is a unit. Here
$v_p(\cross(S))=\sum_{e\ge1}\#\{(k,l):k\in S,\ l\notin S,\ p^e\mid k-l\}$ since $x_k-x_l=2(k-l)$, and
$v(X-P(x_k))\ge-2\max_{n\text{ odd},|n|<|x_k|}v_p(n)\ge-2\sum_{e\ge1}\ind{p^e<|x_k|}$. At level $d=p^e$ the pairs
counted are those inside residue classes, $\sum_c|S\cap C_c|(n_c-|S\cap C_c|)$, and $\ind{p^e<|x_k|}=\frac12w_d(k)$.
Bounding the maximum over $S$ level by level and class by class gives the claim. Levels $p^e\ge K$ contribute nothing:
the classes are singletons and $|x_m|\le h<p^e$.
\end{proof}

\begin{lemma}[Lemma~8A$'$, class count]\label{lem:classcount}
For odd $d$ with $K/3\le d<K$ we have $A_K(d)=3(K-d)$, and $A_K(d)=0$ for $d\ge K$.
\end{lemma}

\begin{proof}
For $d\ge K$ all classes are singletons and all weights vanish. Let $K/3\le d<K$. Then $w_d(m)=0$ if and only if
$|m-J|\le(d-1)/2$, a block of $d$ consecutive indices inside $[0,K)$; so every class contains exactly one atom of weight
$0$, and the others have weight $2$. A class with $n$ atoms contributes $V(n):=\max_t[t(n-t)+2\min(t,n-1)]$ (take the
weight-$0$ atom last), and $V(1)=0$, $V(2)=3$, $V(3)=6$. Since $3d\ge K$, classes have at most three atoms. If
$K/3\le d\le K/2$ there are $K-2d$ classes of size $3$ and $3d-K$ of size $2$, with total $6(K-2d)+3(3d-K)=3(K-d)$. If
$K/2<d<K$ there are $K-d$ pairs and $2d-K$ singletons, with total $3(K-d)$.
\end{proof}

\subsection{The plateau: a bordered Lagrange Gram identity}
In \cite{Chi2Paper} the number of poles equals the size $h$ of the determinant, and the Lagrange basis has the right
dimension. Here $K=h+1$, and the $K$ Lagrange polynomials $f_i:=D/(t-x_i)$ have degree $h$: they span the polynomials of
degree $\le h$, one dimension too many. The adjugate of the $(h+1)$-Gram matrix repairs this.

\begin{lemma}[Lemma~B, bordered identity]\label{lem:bordered}
Let $G^f(X):=(\Lambda_X(f_if_j))_{i,j<K}$. Then
\[
  \Delta_h(X)=(-1)^{K(K-1)/2}\,\frac{\one^{\mathsf T}\adj\bigl(G^f(X)\bigr)\one}{\prod_{i<K}D'(x_i)} .
\]
\end{lemma}

\begin{proof}
Let $H_{h+1}:=(E_{a+b})_{a,b\le h}$. Then $\Delta_h=\adj(H_{h+1})_{hh}$. Let $C$ be the coefficient matrix of the
$f_i$ (row $i$ lists the coefficients of $t^0,\dots,t^h$ in $f_i$). Then $G^f=CH_{h+1}C^{\mathsf T}$, so
$\adj(G^f)=(\det C)^2C^{-\mathsf T}\adj(H_{h+1})C^{-1}$ and $\adj(H_{h+1})=(\det C)^{-2}C^{\mathsf T}\adj(G^f)C$. Column $h$
of $C$ is $\one$, since every $f_i$ is monic of degree $h$; hence $\Delta_h=(\det C)^{-2}\one^{\mathsf T}\adj(G^f)\one$.
Finally $(f_i(x_k))_{i,k}=\operatorname{diag}(D'(x_i))=CW$ with $W=(x_k^a)_{a,k}$, and
$\prod_iD'(x_i)=(-1)^{K(K-1)/2}(\det W)^2$, so $(\det C)^2=(-1)^{K(K-1)/2}\prod_iD'(x_i)$.
\end{proof}

\begin{lemma}[Lemma~9$'$, plateau bound]\label{lem:nine}
For every odd prime $p$, $n_p\le2h\lev(p)+v_p(h!)$.
\end{lemma}

\begin{proof}
Put $N_e(i):=\#\{m<K:m\ne i,\ p^e\mid m-i\}$ and $N(i):=\sum_eN_e(i)=v_p(D'(x_i))$. Note that $N_e(i)\ge1$ needs
$p^e\le K-1=h$. We first prove the \emph{row claim}
\begin{equation}\label{eq:rowclaim}
  v\bigl(G^f_{ij}\bigr)\ge N(i)-2\lev(p)\qquad(i,j<K).
\end{equation}
Fix an odd $t$ and a level $d=p^e$, and let $m_0=m_0(t)$. The class of $i$ modulo $d$ has $N_e(i)+1$ elements in
$[0,K)$, and classes modulo $d$ in $[0,K)$ differ in size by at most one, so the class $m_0$ has at least $N_e(i)$
elements.

\emph{Off-diagonal entries} ($i\ne j$). Here $f_if_j/D=F:=\prod_{m\ne i,j}(t-x_m)$, so $G^f_{ij}=\cM(F)$. Each term of
$F'(t)$ omits one factor. At level $d$ the factors divisible by $d$ are those with $m\equiv m_0$, $m\ne i,j$; after the
omission at least $N_e(i)-2$ remain, whether $m_0\equiv i$ or not. Summing over levels, every term has $v_p\ge
\sum_e(N_e(i)-2\ind{N_e(i)\ge1})\ge N(i)-2\lev(p)$, and \eqref{eq:step3} gives \eqref{eq:rowclaim}.

\emph{Diagonal entries.} Put $g:=\prod_{m\ne i}(t-x_m)$. Then $f_i^2/D=g/(t-x_i)=Q+g(x_i)/(t-x_i)$ with
$Q:=(g(t)-g(x_i))/(t-x_i)$ and $g(x_i)=D'(x_i)$. By Lemma~\ref{lem:pole} the pole part contributes
$-D'(x_i)z^{x_i}(X-P(x_i))$, of valuation at least $N(i)-2\#\{e:p^e<|x_i|\}\ge N(i)-2\lev(p)$, since $|x_i|\le h$. For
$Q$, list the roots $x_m$, $m\ne i$, of $g$ as $\alpha_1,\dots,\alpha_h$ in order of decreasing $v_p(m-i)$; telescoping
gives
\[
  Q(t)=\sum_{r=1}^{h}\prod_{s<r}(x_i-\alpha_s)\prod_{s>r}(t-\alpha_s),
\]
and for every level $e$ the $\alpha_s=x_m$ with $p^e\mid m-i$ form an initial segment of length $N_e(i)$ (note
$x_i-x_m=2(i-m)$). Fix a term $r$ of $Q'(t)$, which omits one more factor of the second product, and a level $d$. If
$r>N_e(i)$, the first product has $N_e(i)$ factors divisible by $d$. If $r\le N_e(i)$ and $m_0\equiv i$, it has $r-1$,
and the remaining $N_e(i)-r$ class members sit at positions $s>r$, at most one of them omitted: in total at least
$N_e(i)-2$. If $r\le N_e(i)$ and $m_0\not\equiv i$, the whole class $m_0$ (at least $N_e(i)$ members, all with
$v_p(m-i)<e$) sits at positions $s>N_e(i)\ge r$: in total at least $N_e(i)-1$. So \eqref{eq:rowclaim} holds for $Q$
too, by \eqref{eq:step3}.

\emph{The minors.} $\one^{\mathsf T}\adj(G^f)\one=\sum_{a,b}\pm M_{ba}$, where $M_{ba}$ is the minor with row $b$ and
column $a$ deleted. By \eqref{eq:rowclaim} and the Leibniz expansion, $v(M_{ba})\ge\sum_{i\ne b}(N(i)-2\lev(p))\ge
\sum_iN(i)-2h\lev(p)-\max_bN(b)$. Lemma~\ref{lem:bordered} and $v_p(\prod_iD'(x_i))=\sum_iN(i)$ give
$v(\Delta_h)\ge-2h\lev(p)-\max_bN(b)$. Finally $N(b)=v_p(b!\,(h-b)!)\le v_p(h!)$.
\end{proof}

\begin{remark}
On the exact profiles computed for $h=16,\dots,256$ the bound of Lemma~\ref{lem:nine} is never violated, and it is
attained at most single-level plateau primes; for example at $h=128$ the primes $13,17,19,23,31,37,41,43$ have
$n_p=256+v_p(128!)$ exactly. The row claim \eqref{eq:rowclaim} was checked entrywise at $h=16,24,32,40$ and $p\le13$, with
minimum slack $0$.
\end{remark}

\subsection{The sum over the odd primes}

\begin{theorem}[Theorem~A$'$]\label{thm:A}
For every $\eps>0$ and all large even $h$, $\sum_{p\text{ odd}}n_p\log p\le(\frac43+\eps)h^2$.
\end{theorem}

\begin{proof}
For $p<K/3$ use Lemma~\ref{lem:nine}:
\begin{align*}
  \sum_{p<K/3}\bigl(2h\lev(p)+v_p(h!)\bigr)\log p&\le2h\bigl(\theta(K/3)+\psi(h)-\theta(h)\bigr)+\log h!\\
  &=\tfrac23h^2+o(h^2)
\end{align*}
by the prime number theorem. For $K/3\le p<K$ we have $p^2\ge K^2/9\ge K$ for $h\ge8$, so only the level $e=1$ occurs,
and Lemmas~\ref{lem:eight} and~\ref{lem:classcount} give $n_p\le3(K-p)$. By the prime number theorem and partial
summation, $\sum_{K/3\le p<K}3(K-p)\log p=3h^2\int_{1/3}^1(1-x)\,dx+o(h^2)=\frac23h^2+o(h^2)$. For $p\ge K$,
$n_p\le0$ by Lemma~\ref{lem:eight}.
\end{proof}

The finite sums of these bounds, divided by $h^2$, are $1.3380$, $1.3350$ and $1.3337$ at $h=1442$, $20836$ and
$875032$. The prime $5$ needs no separate treatment: $z$ and $1-y=4z$ are units at $\fp=(\sqrt5)$, so every step above
holds there with $v(5)=1$.

% ===========================================================================
\section{Denominators II: the prime 2}\label{sec:two}
% ===========================================================================

Let $v$ be the $2$-adic valuation of $\KK$; $2$ is inert, $v(2)=1$, $v(z)=0$ and $v(1-y)=2$.

\emph{Rescaling.} Put $\tau:=t/2$. The atoms become $\tau=s+\frac12$ ($s\ge0$), the poles $\xi_k:=k-J\in\ZZ$, and
$D(t)=2^K\tilde D(\tau)$ with $\tilde D(\tau):=\prod_{k<K}(\tau-\xi_k)$. With $\tilde\cM(F):=\sum_{s\ge0}z^{2s+1}F'(s+\frac12)$
and $\tilde\Lambda(g):=\tilde\cM(g/\tilde D)$ we get $\Lambda(t^m)=2^{m-K-1}\tilde\Lambda(\tau^m)$, hence
$\Delta_h=2^{-3h}\tilde\Delta_h$ with $\tilde\Delta_h:=\det(\tilde\Lambda_X(\tau^{i+j}))$, since
$\sum_{i<h}2i-h(K+1)=-3h$.

\emph{Anchored basis.} Put $a_i:=\prod_{l<i}(\tau-\xi_l)$ and $b_j:=\prod_{l<j}(\tau-\xi_{K-1-l})$. Both families are
monic of degrees $i$ and $j$, so $\det(G_{ij})=\tilde\Delta_h$ for $G_{ij}:=\tilde\Lambda_X(a_ib_j)$.

\begin{lemma}[Lemma~10$'$]\label{lem:ten}
Let $i,j<h$.
\begin{enumerate}
\item[(L)] If $\pi:=i+j-K\ge1$, then $v(G_{ij})\ge-\pi-\lfloor\log_2\pi\rfloor-2$. If $\pi=0$, then $G_{ij}=0$.
\item[(U)] If $N_0:=K-1-i-j\ge1$, then $v(G_{ij})\ge N_0-\lfloor\log_2N_0\rfloor$. If $N_0=0$, then $v(G_{ij})=2$.
\end{enumerate}
Consequently $n_2\le h\lfloor\log_2h\rfloor+3h$.
\end{lemma}

\begin{proof}
(L) The products $a_i$ and $b_j$ cover the index ranges $[0,i-1]$ and $[K-j,K-1]$, which overlap in $\pi$ indices, so
$a_ib_j/\tilde D=F:=\prod_{l=K-j}^{i-1}(\tau-\xi_l)=2^{-\pi}H(2\tau)$ with $H(U):=\prod_l(U-2\xi_l)\in\ZZ[U]$. So
$F'(s+\frac12)=2^{1-\pi}H'(2s+1)$, and by \eqref{eq:mahler} (with $s$ in place of $(t-1)/2$),
$\tilde\cM(F)=z\sum_{b<\pi}\Delta^bG(0)y^b/(1-y)^{b+1}$, $G(s)=2^{1-\pi}H'(2s+1)$. The Newton coefficients of
$s\mapsto(2s+1)^k$ are $\Delta^b[(2s+1)^k](0)=\sum_{j\ge b}\binom kj2^j\,b!\,S(j,b)$, of valuation $\ge b+v_2(b!)$, while
$y^b/(1-y)^{b+1}$ has valuation $-2(b+1)$. Hence
$v(G_{ij})\ge1-\pi+\min_{b<\pi}(-2-s_2(b))\ge-\pi-\lfloor\log_2\pi\rfloor-2$. There is no $X$-part. If $\pi=0$ then
$F=1$ and $\tilde\cM(1)=0$.

(U) Now the ranges leave a gap of $N_0+1$ indices, and with $\zeta_l:=\xi_i+l$,
\[
  \frac{a_ib_j}{\tilde D}=\prod_{l=0}^{N_0}\frac1{\tau-\zeta_l}=\frac1{N_0!}\sum_{l=0}^{N_0}(-1)^{N_0-l}\binom{N_0}l\frac1{\tau-\zeta_l}.
\]
For an integer $\zeta$, $\tilde\cM(1/(\tau-\zeta))=-4\sum_sz^{2s+1}/(2s+1-2\zeta)^2=-4z^{2\zeta}(X-S(\zeta))$ (at $X=T$),
where $S(\zeta+1)-S(\zeta)=-z^{-1-2\zeta}/(2\zeta+1)^2$; so $S(\zeta)$ has only odd denominators and $v(S(\zeta))\ge0$.
With $z^{2\zeta_l}=z^{2\xi_i}y^l$ and $S(\zeta_l)=S(\xi_i)-\sum_{r<l}z^{-1-2\xi_i-2r}g(r)$, $g(r):=(s_0+2r)^{-2}$,
$s_0:=2\xi_i+1$,
\[
  G_{ij}=-\frac4{N_0!}\Bigl[z^{2\xi_i}(y-1)^{N_0}\bigl(X-S(\xi_i)\bigr)+z^{-1}\sum_l(-1)^{N_0-l}\binom{N_0}l\sum_{r<l}y^{l-r}g(r)\Bigr].
\]
The first part has valuation $2N_0$, since $y-1=-4z$. In the second, write $y^{l-r}=y\sum_e\binom{l-r-1}e\eta^e$ with
$\eta:=y-1$ ($v(\eta)=2$) and $g(r)=\sum_m\gamma_m\binom rm$; expanding
$g(r)=s_0^{-2}\sum_k(-1)^k(k+1)(2/s_0)^kr^k$ (convergent in $\ZZ_2$) and using $\Delta^mr^k(0)=m!\,S(k,m)$ gives
$v_2(\gamma_m)\ge m+v_2(m!)=2m-s_2(m)$. The identities $\sum_{r<l}\binom rm\binom{l-1-r}e=\binom l{m+e+1}$ and
$\sum_l(-1)^{N_0-l}\binom{N_0}l\binom lM=\ind{M=N_0}$ collapse the second part to $z^{-1}y\sum_{m+e=N_0-1}\gamma_m\eta^e$,
of valuation $\ge2N_0-3-\lfloor\log_2N_0\rfloor$. With $v_2(N_0!)=N_0-s_2(N_0)$ this gives
$v(G_{ij})\ge N_0-1+s_2(N_0)-\lfloor\log_2N_0\rfloor\ge N_0-\lfloor\log_2N_0\rfloor$; the $X$-part has valuation
$N_0+2+s_2(N_0)$. If $N_0=0$ then $G_{ij}=-4z^{2\xi_i}(X-S(\xi_i))$, of valuation exactly $2$.

\emph{Assignment.} In all cases $v(G_{ij})\ge K-i-j-\lfloor\log_2h\rfloor-2$, since $\pi,N_0\le h$. By the Leibniz
expansion, $v(\tilde\Delta_h)\ge\sum_{i<h}(K-i-\sigma(i))-h(\lfloor\log_2h\rfloor+2)=-h\lfloor\log_2h\rfloor$ for every
permutation $\sigma$, and $n_2=-v(\Delta_h)=3h-v(\tilde\Delta_h)$.
\end{proof}

The mechanism is $v_2(1-y)=2$: both triangles (L) and (U) then have slope $1$ in $i+j$, and the permutation sum of the
bound is constant. In \cite{Chi2Paper} the corresponding quantity is $v_2(1-q^{-2})=3$, which costs $\frac14\log2$ per
$h^2$. Section~\ref{sec:M1} shows what happens when it is $1$. The exact values $n_2=54,96,142,240,350,576,830$ at
$h=16,24,32,48,64,96,128$ lie below the bound of Lemma~\ref{lem:ten} ($112,168,256,384,576,864,1280$), and
$n_2/(h\log_2h)$ rises from $0.84$ to $0.94$ for $h\le256$.

% ===========================================================================
\section{The size of \texorpdfstring{$\Delta_h(T)$}{Delta_h(T)}: Theorem B\texorpdfstring{$'$}{'}}\label{sec:analytic}
% ===========================================================================

\subsection{The equilibrium problem and weak duality}
Put
\[
  \Psi(x):=\tfrac12\int_{x-1}^{x+1}\log|s|\,ds,\qquad v(x):=x\log z-\Psi(x).
\]
$\Psi$ is even, increasing on $[0,\infty)$, and $\Psi(0)=-1$. The function $v$ is continuous, and
$v'(x)=\log z-\frac12\log|(x+1)/(x-1)|$ for $x\ne\pm1$. For $x>0$, $x\ne1$, the logarithm is positive, so
$v'<\log z<0$; for $-1<x\le0$, $v'(x)=\log z+\artanh|x|<0$ if and only if $|x|<\tanh(\log(1/z))=0.8945\ldots$. Hence
\begin{equation}\label{eq:vmono}
  v\text{ is strictly decreasing on }(-0.894,\infty)\supset[-1/h,\infty)\qquad(h\ge2).
\end{equation}
For a finite positive measure $\rho$ with compact support put $U^\rho(x):=\int\log|x-s|\,d\rho(s)$ and
$E(\rho):=\iint\log|x-s|\,d\rho(x)\,d\rho(s)$.

Call $\mu$ \emph{admissible} if it is a density with $0\le\mu\le1$, mass $1$ and $\operatorname{supp}\mu\subset[0,2]$.
For admissible $\mu$ and $\lambda\in\RR$ put
\[
  F_\mu:=2U^\mu+v,\qquad \Phi(\mu,\lambda):=-E(\mu)+\lambda+\int_0^\infty(F_\mu-\lambda)_+\,dx .
\]
If $\rho$ is a density with $0\le\rho\le1$, mass $1$ and compact support, then $E(\rho-\mu)\le0$
\cite[Lemma~I.1.8]{SaffTotik}, i.e.\ $E(\rho)\le2\int U^\mu\,d\rho-E(\mu)$, so
\begin{equation}\label{eq:weakdual}
  E(\rho)+\int v\,d\rho\le-E(\mu)+\lambda+\int(F_\mu-\lambda)\,d\rho .
\end{equation}
The supremum of the left side over such $\rho$ supported in $[0,\infty)$ is the constrained equilibrium constant $C$;
by \eqref{eq:weakdual}, $C\le\Phi(\mu,\lambda)$.

\subsection{The theorem}

\begin{theorem}[Theorem~B$'$]\label{thm:B}
Let $\mu$ be admissible and $\lambda\ge-2.96317$, and suppose that $F_\mu-\lambda\le-1.13-0.94(x-4)$ on $[4,\infty)$.
Then for every even $h\ge100$,
\[
  \log|\Delta_h(T)|\le h^2\,\Phi(\mu,\lambda)+E_{\mathrm{an}}(h),
\]
where $c_1(h):=\log(2h)+1+\log(2+\log(2h+1))$ and
\[
  E_{\mathrm{an}}(h):=2h\log(4h)+h\log4+(2h+1)\log3+h\,c_1(h)+7.742h-h\log h+1 .
\]
In particular $\limsup_{h\to\infty}h^{-2}\log|\Delta_h(T)|\le\Phi(\mu,\lambda)$.
\end{theorem}

The proof follows the proof of Theorem~B in \cite{Chi2Paper} in the first two steps and then departs from it: the poles
now lie on both sides of the atoms, which is handled by an exact Jensen-type bound (step~4), and the discretisation
spreads the particles into cells (step~5) instead of using a lattice energy inequality.

\begin{proof}
Write $\phi:=1/D$.

\emph{1. Node expansion.} For $a,b\in\QQ[t]$, $(ab\phi)'=a'b\phi+ab'\phi+ab\phi'$, so
$\Lambda(ab)=\sum_{t\text{ odd}\ge1}(a,a')(t)\,M_t\,(b,b')(t)^{\mathsf T}$ with
$M_t:=z^t\bigl(\begin{smallmatrix}\phi'&\phi\\\phi&0\end{smallmatrix}\bigr)(t)$. Let $H^{(t_{\max})}$ be the Hankel matrix of
the sums truncated at $t\le t_{\max}$. Each node $t$ carries a \emph{value slot} (the row $(t^i)_{i<h}$) and a
\emph{derivative slot} (the row $(it^{i-1})_{i<h}$). By the Cauchy--Binet formula, $\det H^{(t_{\max})}$ is a sum over
\emph{patterns} $n=(n_t)$, $n_t\in\{0,1,2\}$, $\sum_tn_t=h$, and over pairs $(S,S')$ of slot sets with $n_t$ slots at
node $t$, of $\pm\det E_S\,\det\mathbb M_{S,S'}\,\det E_{S'}$, where $E_S$ is the matrix of the rows in $S$ and
$\mathbb M$ is block diagonal with blocks $M_t$. There are at most $4^h$ pairs $(S,S')$ per pattern.

\emph{2. Slot weights.} $|\det\mathbb M_{S,S'}|\le\prod_t(z^tm_t)^{n_t}$ with $m_t:=\max(|\phi(t)|,|\phi'(t)|)$
(for $n_t=2$ note $|\det M_t|=z^{2t}\phi(t)^2$). The $K$ numbers $t-x_m$ are distinct odd integers, so
$|\phi'/\phi|(t)\le\sum_m1/|t-x_m|\le2\sum_{k=0}^h1/(2k+1)\le2+\log(2h+1)$, and $m_t\le|\phi(t)|(2+\log(2h+1))$.

\emph{3. Vandermonde determinants with derivative rows.} $|\det E_S|\le(4h)^h\prod_{t<t'}|t-t'|^{n_tn_{t'}}$. Indeed,
replace the derivative-only slots $s_1,\dots,s_r$ of $S$ by value rows at free variables; the result is a confluent
Vandermonde determinant $\pm\prod_{a<b}(a-b)^{n_an_b}$, and $\det E_S$ is obtained by applying
$\partial_{s_1}\cdots\partial_{s_r}$ and evaluating at the nodes. By induction, after $k$ derivatives the expression is a sum of
at most $(h-1)^k$ terms, each a product of $k$ exponents (each $\le2$) times $\prod(a-b)^{e'_{ab}}$ with
$0\le e'_{ab}\le n_an_b$. At the nodes $|a-b|\ge2>1$, so each term is at most $2^k\prod|a-b|^{n_an_b}$, and
$(2(h-1))^r\le(4h)^h$.

\emph{4. The pole product.} For odd $t$ the numbers $t-x_m$ are the odd integers in $[t-h,t+h]$. For every odd $u$,
$\log|u|\ge\frac12\int_{u-1}^{u+1}\log|s|\,ds$, by concavity of $\log$ on $(0,\infty)$ and on $(-\infty,0)$ (the
singular endpoint is integrable). The intervals $[u-1,u+1]$ tile $[t-h-1,t+h+1]$, so with $x:=t/h$ and $s=u/h$,
\begin{align*}
  \log|D(t)|&\ge\tfrac12\int_{t-h-1}^{t+h+1}\log|u|\,du\\
  &=(h+1)\log h+h\Psi(x)+\tfrac h2\Bigl[\int_{x+1}^{x+1+1/h}+\int_{x-1-1/h}^{x-1}\Bigr]\log|s|\,ds .
\end{align*}
Each bracketed integral is over an interval of length $1/h$, hence at least $\int_{-1/(2h)}^{1/(2h)}\log|s|\,ds=-(\log(2h)+1)/h$.
So $\log|D(t)|\ge(h+1)\log h+h\Psi(x)-\log(2h)-1$ and, with step~2,
\begin{equation}\label{eq:nodeweight}
  t\log z+\log m_t\le h\,v(t/h)-(h+1)\log h+c_1(h).
\end{equation}
By steps 1--4, the pattern $n$ contributes at most $4^h(4h)^{2h}e^{D(n)}$ with
$D(n):=\sum_{t\ne t'}n_tn_{t'}\log|t-t'|+\sum_tn_t\bigl(h\,v(t/h)-(h+1)\log h+c_1(h)\bigr)$.

\emph{5. Discrete to continuum.} Spread the $n_t$ slots of node $t$ uniformly over $I_t:=[(t-2)/h,t/h]$: the density
$\rho_n:=\sum_t\frac{n_t}2\one_{I_t}$ satisfies $0\le\rho_n\le1$, has mass $1$, and is supported in $[-1/h,\infty)$. By
\eqref{eq:vmono}, $v(t/h)\le v$ on $I_t$, so $\sum_tn_th\,v(t/h)\le h^2\int v\,d\rho_n$. For $t\ne t'$ with
$|t-t'|=2d$, writing $x=(t-2+2a)/h$, $a\in[0,1]$,
\begin{gather*}
  h^2\iint_{I_t\times I_{t'}}\log|x-s|\,d\rho_n\,d\rho_n=n_tn_{t'}\bigl(\log|t-t'|-\log h-\operatorname{gap}(d)\bigr),\\
  \operatorname{gap}(d):=\log d-\int_0^1\!\!\int_0^1\log|d+a-a'|\,da\,da',
\end{gather*}
and $0\le\operatorname{gap}(d)\le0.1137/d^2$ (with equality at $d=1$ up to rounding). A cell with itself gives
$n_t^2(\log2-\log h-\frac32)$. Summing, and using $n_{t'}\le2$ and $\sum_tn_t^2\le2h$,
\[
  \sum_{t\ne t'}n_tn_{t'}\log|t-t'|\le h^2E(\rho_n)+h^2\log h+0.748h+1.614h
\]
($4\cdot0.1137\,\zeta(2)\le0.748$ and $2(\frac32-\log2)\le1.614$). The terms in $\log h$ combine to
$h^2\log h-h(h+1)\log h=-h\log h$, so
\[
  D(n)\le h^2J(\rho_n)-h\log h+h\,c_1(h)+2.362h,\qquad J(\rho):=E(\rho)+\int v\,d\rho .
\]

\emph{6. Dual, the cell below $0$, and far nodes.} By \eqref{eq:weakdual},
$h^2J(\rho_n)\le h^2(-E(\mu)+\lambda)+h^2\int(F_\mu-\lambda)\,d\rho_n$. On $[0,4]$ the integral is at most
$\int_0^4(F_\mu-\lambda)_+dx$, since $\rho_n\le1$. On $[-1/h,0]$, which only the cell $I_1$ meets, we have
$U^\mu(x)\le\log(2+1/h)$ (as $\operatorname{supp}\mu\subset[0,2]$) and $v(x)\le|\log z|/h-\Psi(0)=1+1.4437/h$; hence
$F_\mu\le2.4108$ and $F_\mu-\lambda\le2.4108+2.96317\le5.374<5.38$ for $h\ge100$, and the contribution to $h^2J$ is at
most $5.38h$. On $[4,\infty)$, by hypothesis, a slot at a node $t\ge4h+3$ contributes at most $-a_t$ with
$a_t:=h(1.13+0.94((t-2)/h-4))\ge1.13h$; the node $4h+1$, whose cell straddles $4$, has $F_\mu-\lambda<0$ on its part
beyond $4$. So for every pattern
\[
  h^2J(\rho_n)\le h^2\Phi(\mu,\lambda)+5.38h-\sum_{t\ge4h+3}n_ta_t .
\]
The $2h+1$ odd nodes $t\le4h+1$ carry at most $3^{2h+1}$ patterns, and the nodes beyond contribute at most
$\prod_{t\ge4h+3}(1+e^{-a_t}+e^{-2a_t})\le\exp\bigl(2e^{-1.13h}/(1-e^{-1.88})\bigr)\le e$. Collecting,
\begin{multline*}
  \log|\det H^{(t_{\max})}|\le h\log4+2h\log(4h)+(2h+1)\log3\\
  +h^2\Phi(\mu,\lambda)+5.38h-h\log h+h\,c_1(h)+2.362h+1,
\end{multline*}
uniformly in $t_{\max}$. As $t_{\max}\to\infty$ the truncated matrices converge entrywise to $(\Lambda(t^{i+j}))$, and
the determinant is continuous.
\end{proof}

\begin{remark}
For every even $h\ge2$ the same bound holds with $7.742$ replaced by $9.592$: for any $h$, the bound
$U^\mu(x)\le\int_1^2\log(s+1)\,ds=0.9095\ldots$ for $x\in[-1/h,0]$ (because $\mu\le1$ and $\operatorname{supp}\mu\subset[0,2]$)
gives $F_\mu-\lambda\le7.23$ on $[-1/h,0]$. Only the $O(h)$ term changes.
\end{remark}

\subsection{The certificate}\label{sec:cert}

\begin{theorem}\label{thm:cert}
There is an admissible pair $(\mu,\lambda)$ satisfying the hypotheses of Theorem~\ref{thm:B} with
\[
  \Phi(\mu,\lambda)\le-1.3592161834586374 .
\]
\end{theorem}

The proof is a finite computation in ball arithmetic (Arb \cite{Arb}, through python-flint \cite{pythonflint}) at $128$
bits. The density $\mu$ is constant on $N$ equal cells of $[0,2]$, with rational heights in $[0,1]$ and mass exactly $1$
(checked in exact rational arithmetic). The energy $E(\mu)$ is computed from the exact double antiderivative
$\frac12u^2\log|u|-\frac34u^2$ as a Toeplitz sum over cell offsets. The grid has step $d/r$ ($d=2/N$ the cell width,
$r$ even, so cell midpoints are grid points), and on each grid interval $I=[a,b]$:
\begin{itemize}
\item the three cells nearest to $I$ are treated exactly: each cell potential is unimodal with its minimum at the cell
  midpoint, so its maximum on $I$ is at an endpoint that can be chosen exactly;
\item the far field $f=2\sum_{\text{far}}\mu_jU_j$ (every far cell is at distance $\ge d$ from $I$) satisfies
  $f''\ge-4/d$, because $f''=-2\int_{\text{far}}\mu(s)(x-s)^{-2}ds$ and $\mu\le1$; and $v''(x)=1/(x^2-1)$ is $\ge-1/(1-b^2)$
  for $b<1$ and $\ge0$ for $a>1$. A function $g$ with $g''\ge-\kappa$ on $I$ satisfies
  $\sup_Ig\le\max(g(a),g(b))+\kappa(b-a)^2/8$; this is applied to $f+v$ with $\kappa=4/d+1/(1-b^2)$ (if $b<1$) or
  $\kappa=4/d$ (if $a>1$). On the two intervals touching $x=1$, $v$ is bounded by $v(a)$;
\item $\int(F_\mu-\lambda)_+\le\sum_I|I|\max(S_I-\lambda,0)$, with $S_I$ the resulting upper bound on $I$; on
  $[4,\infty)$, $F_\mu(x)\le2\log x+x\log z-\log(x-1)$, which is $<\lambda$ and decreasing there.
\end{itemize}
A first version with per-cell maxima everywhere loses the cancellation between neighbouring cells and gives only
$-1.313$ at $N=100$. Table~\ref{tab:cert} lists the results; the last row proves Theorem~\ref{thm:cert}.

\begin{table}[ht]
\centering\small
\begin{tabular}{@{}rrcccc@{}}
\toprule
$N$ & $r$ & $-E(\mu)$ & $\lambda$ & $\int(F_\mu-\lambda)_+\le$ & certified $\Phi$\\
\midrule
200 & 16 & 1.21902586 & $-2.96219282$ & 0.38673 & $-1.35643705$\\
400 & 32 & 1.21902867 & $-2.96297388$ & 0.38528 & $-1.35866119$\\
800 & 32 & 1.21902407 & $-2.96310589$ & 0.38505 & $-1.35903096$\\
1600 & 32 & 1.21902464 & $-2.96316943$ & 0.384929 & $-1.35921618345864$\\
\bottomrule
\end{tabular}
\caption{Dual certificates for Theorem~\ref{thm:B} (decimal values rounded; $\lambda$ is a rational number). For
$N=1600$ the certified Arb ball is $[-1.35921618345864\pm2.6\cdot10^{-15}]$, with upper endpoint
$-1.3592161834586374$.}\label{tab:cert}
\end{table}

The side conditions of Theorem~\ref{thm:B} hold for the certified pair: $\lambda=-2.96316943\ldots\ge-2.96317$; at
$x=4$, $2\log x+x\log z-\log(x-1)-\lambda=-1.137\ldots$, and the slope of this bound on $[4,\infty)$ is at most
$2/x+\log z\le-0.94$.

\begin{remark}[the equilibrium]\label{rem:equil}
An active-set quadratic program (and two independent codes) give the primal value $C\approx-1.3594016$, with maximiser
saturated ($\rho=1$) on $[0,0.5475]$, a free band up to $1.9525$, and $\lambda^*\approx-2.96317$; so the certificate is
within about $1.9\cdot10^{-4}$ of the optimum. The finite-$h$ values $h^{-2}\log|\Delta_h(T)|$, computed in ball
arithmetic, are $-1.4379$, $-1.4195$, $-1.4093$, $-1.4015$, $-1.3930$, $-1.3876$, $-1.3815$, $-1.3781$ at $h=16$, $32$,
$48$, $64$, $96$, $128$, $192$, $256$. They increase towards the constant from below; fits of the form $a+b\log h/h\,(+c/h)$
give limits between $-1.373$ and $-1.364$, below $C$, so these data do not indicate that Theorem~\ref{thm:B} is sharp.
This costs nothing, since only the upper bound is used.
\end{remark}

% ===========================================================================
\section{Nonvanishing: Lemma P\texorpdfstring{$'$}{'}}\label{sec:nonvan}
% ===========================================================================

\subsection{Setting}
Let $\ell\ge7$ be prime, $h:=3\ell-1$ (even) and $K=3\ell$. The classes modulo $\ell$ are $C_\gamma=\{\gamma,\gamma+\ell,
\gamma+2\ell\}$, $\gamma<\ell$, all of size $3$. Since $|m-i|<3\ell<\ell^2$ there is a single level, and
$N(i)=v_\ell(D'(x_i))=2$ for every $i$. Let $\fl$ be a prime of $\KK$ above $\ell$, $v:=v_\fl$, and $\zeta:=z^\ell\bmod\fl$;
then $\zeta=z$ if $\ell$ splits ($\ell\equiv\pm1\bmod5$) and $\zeta=\sigma(z)=-1/z$ if $\ell$ is inert. Since $\ell\nmid10$,
the elements $z$, $y$, $1-y=4z$ and $1-\zeta^{\pm2}$ are $\fl$-units.

\emph{Positions.} $x_{\gamma+r\ell}=x_\gamma+2r\ell$ with $x_\gamma=2\gamma-3\ell+1\in[-3\ell+1,-\ell-1]$: in each class the
pole with $r=0$ is large and negative ($|x|>\ell$), the pole with $r=1$ is small ($|x|<\ell$), and the pole with $r=2$ is
large and positive. For odd $t$ we have $\ell\mid t-x_m$ exactly for $m$ in one class $C_{\gamma(t)}$, and the odd $t\ge1$
with $\gamma(t)=\gamma$ form the progression $t=x_\gamma+\ell Y$, $Y\in\{3,5,7,\dots\}$. Put $t_\gamma:=x_\gamma+3\ell$.

\begin{lemma}[Lemma~P$'$]\label{lem:P}
Let $c_k:=[X^k]\Delta_h$. Then $v_\ell(c_0)=-2K$ and $v_\ell(c_k)\ge-2K+2k$ for all $k$. Hence $\Delta_h(a/b)\ne0$ for every
$a/b\in\QQ$ with $v_\ell(b)\le1$.
\end{lemma}

\begin{proof}
\emph{1. Lemma B.} $\ell^{2K}\Delta_h(X)=(-1)^{K(K-1)/2}u_\ell\,\one^{\mathsf T}\adj(G^f(X))\one$ with
$u_\ell:=\ell^{2K}/\prod_iD'(x_i)$ an $\ell$-unit.

\emph{2. The $X$-part.} $B(f_if_j)=0$ for $i\ne j$, because $f_if_j/D$ is a polynomial, and $B(f_i^2)=-D'(x_i)z^{x_i}$, of
valuation exactly $2$. So $G^f(X)=G^f(0)+X\operatorname{diag}(d_i)$ with $v(d_i)=2$.

\emph{3. Integrality and block diagonality.} By the row claim \eqref{eq:rowclaim} with $N(i)-2\lev(\ell)=0$, $G^f(0)$ is
$\fl$-integral. On each progression, $\cM$ restricted to $\gamma$ is
$z^{t_\gamma}\sum_k z^{2\ell k}F'(t_\gamma+2\ell k)=z^{t_\gamma}\sum_b\Delta^bG_\gamma(0)z^{2\ell b}/(1-z^{2\ell})^{b+1}$,
$G_\gamma(k):=F'(t_\gamma+2\ell k)$, with $1-z^{2\ell}\equiv1-\zeta^2$ a unit; so \eqref{eq:step3} holds progression by
progression. Off the diagonal, $F=\prod_{m\ne i,j}(t-x_m)$ has on the progression $\gamma$ exactly
$3-\ind{\gamma=c(i)}-\ind{\gamma=c(j)}$ factors divisible by $\ell$ ($c(i)$ is the class of $i$), and $F'$ keeps all but
one; so a unit is possible only if $c(i)=c(j)=\gamma$. On the diagonal, the polynomial part on a progression
$\gamma\ne c(i)$ has valuation at least $3-1=2$: the class $\gamma$ has three members, all after the initial segment of
the telescoping order, and the derivative loses at most one. Hence $G^f(0)\bmod\fl$ is block diagonal over the
classes, each block receiving only its own progression and its own pole terms.

\emph{4. The blocks.} Fix $\gamma$ with members $m_r=\gamma+r\ell$, put $W(t):=\prod_{m\notin C_\gamma}(t-x_m)$ and
$U_\gamma:=\prod_{m\notin C_\gamma}(x_\gamma-x_m)$, so $W(t)\equiv U_\gamma\bmod\ell$ on the progression. Off the diagonal
($r\ne s$, $q$ the third index) $F=W(t)(t-x_{m_q})$ and $F'\equiv W\equiv U_\gamma$. For the diagonal polynomial part,
$Q=W\cdot(t+x_{m_r}-x_{m_a}-x_{m_b})+p(x_{m_r})\,W[x_{m_r},t]$ with $\{a,b\}$ the other indices, $p(t)=(t-x_{m_a})(t-x_{m_b})$ and
$p(x_{m_r})\equiv0\bmod\ell^2$; so $Q'\equiv U_\gamma$ too. In the progression sum, $\Delta^0G_\gamma(0)\equiv U_\gamma$ and
$\Delta^bG_\gamma(0)\equiv0$ for $b\ge1$, so it is $\equiv U_\gamma z^{t_\gamma}/(1-z^{2\ell})\equiv U_\gamma z^{x_\gamma}\zeta^3/(1-\zeta^2)$.
The pole terms at $X=0$ are $D'(x)z^xP(x)$ with $\ell^2P(x_{m_r})\equiv\zeta,0,-\zeta^{-1}$ for $r=0,1,2$ (the only odd
multiples of $\ell$ below $|x|\le3\ell-1$ in absolute value are $\pm\ell$), $D'(x_{m_r})/\ell^2\equiv U_\gamma\cdot(8,-4,8)_r$,
and $z^{x_{m_r}}\equiv z^{x_\gamma}\zeta^{2r}$. Hence
\begin{gather*}
  G^f(0)|_{C_\gamma}\equiv e_\gamma M,\qquad e_\gamma:=U_\gamma z^{x_\gamma},\\
  M:=c\,\one\one^{\mathsf T}+\operatorname{diag}(8\zeta,0,-8\zeta^3),\qquad c:=\frac{\zeta^3}{1-\zeta^2}.
\end{gather*}

\emph{5. The bordered determinant.} For $M=\operatorname{diag}(d_0,0,d_2)+c\one\one^{\mathsf T}$ we have
$\det M=c\,d_0d_2=-64c\zeta^4$ and $\one^{\mathsf T}\adj(M)\one=d_0d_2=-64\zeta^4$, both units. For the block-diagonal
reduction,
\[
  \one^{\mathsf T}\adj\Bigl(\bigoplus_\gamma e_\gamma M\Bigr)\one=(\det M)^{\ell-1}\,\one^{\mathsf T}\adj(M)\one\,
  \Bigl(\prod_\gamma e_\gamma^3\Bigr)\sum_\gamma e_\gamma^{-1}.
\]
Modulo $\ell$, $U_\gamma\equiv\prod_{m\not\equiv\gamma}2(\gamma-m)\equiv2^{3(\ell-1)}((\ell-1)!)^3\equiv-1$ for every $\gamma$
(Fermat, Wilson). So $\sum_\gamma e_\gamma^{-1}\equiv-z^{3\ell-1}\sum_{\gamma<\ell}z^{-2\gamma}=-z^{3\ell-1}(1-z^{-2\ell})/(1-z^{-2})$, and
$1-z^{-2\ell}\equiv1-\zeta^{-2}$ and $1-z^{-2}=-4/z$ are units. Hence $\one^{\mathsf T}\adj(G^f(0))\one$ is an $\fl$-unit, and
$v_\ell(c_0)=-2K$ by step~1.

\emph{6. The gap.} $[X^k]\one^{\mathsf T}\adj(G^f(X))\one$ is a sum of products containing $k$ factors $d_i$, so it has
valuation $\ge2k$. If $v_\ell(b)\le1$, the $k$-th term of $\Delta_h(a/b)$ has valuation $\ge-2K+2k-k>-2K$ for $k\ge1$.
\end{proof}

Every prime $\ell\ge7$ works, split or inert. In \cite{Chi2Paper} the partial sums reach past $3\ell$ and the block
constants are fixed rationals that must avoid $\ell$, so finitely many primes are excluded; here only $n=\pm\ell$ enters.
Steps~4--5 were checked entrywise modulo $\ell$ for $\ell=7,11,\dots,41$ (split and inert), and the conclusion on the exact
$\Delta_h$ for $\ell=7,\dots,29$.

% ===========================================================================
\section{Proof of Theorem~\ref{thm:main}}\label{sec:proof}
% ===========================================================================

Suppose that $T=a/b$ with $a\in\ZZ$, $b\ge1$. Let $\mathcal H:=\{3\ell-1:\ell\ge37\text{ prime},\ \ell^2\nmid b\}$, an
infinite set of even integers $h\ge110$. For $h\in\mathcal H$, Lemma~\ref{lem:P} gives $\Delta_h\not\equiv0$ and
$F_h(a/b)\ne0$, where $F_h\in\ZZ[X]$ is the primitive polynomial of Section~\ref{sec:construction}, of degree $\le h$. We
have $\log|F_h(T)|=\log|\Delta_h(T)|+\sum_pn_p\log p$, and along $\mathcal H$:
\begin{itemize}
\item odd primes: $\limsup h^{-2}\sum_{p\text{ odd}}n_p\log p\le\frac43$ (Theorem~\ref{thm:A});
\item $p=2$: $n_2=O(h\log h)$ (Lemma~\ref{lem:ten});
\item size: $\limsup h^{-2}\log|\Delta_h(T)|\le-1.3592161834586374$ (Theorems~\ref{thm:B} and~\ref{thm:cert}).
\end{itemize}
So $\limsup_{h\in\mathcal H}h^{-2}\log|F_h(T)|\le-0.0258828\ldots<0$, contradicting Proposition~\ref{prop:criterion}.
Hence $T=\chi_2(\sqrt5-2)$ is irrational, and so is $24T=\pi^2-18\log^2\varphi$ by \eqref{eq:closedform}. \qed

\begin{remark}[effectivity]\label{rem:effective}
Every bound above is explicit at finite $h$ except the prime number theorem in Theorem~\ref{thm:A}. At a single $h$ one
can instead evaluate the proved finite-$h$ bound
\[
  \log|F_h(T)|\le h^2\Phi+E_{\mathrm{an}}(h)+B_{\mathrm{odd}}(h)+h(\lfloor\log_2h\rfloor+3)\log2,
\]
where $B_{\mathrm{odd}}(h)=\sum_{3\le p<K/3}(2h\lev(p)+v_p(h!))\log p+\sum_{K/3\le p<K}\min(2h+v_p(h!),3(K-p))\log p$ is an
exact finite sum. In exact arithmetic this bound is negative at every even $h$ with $2432\le h\le40000$, and at every
sampled $h$ up to $8.75\cdot10^5$, with $|\text{bound}|/h^2$ increasing towards $0.0259$. A proof that it stays negative
for all larger $h$ would need explicit error terms in the prime number theorem, which we have not carried out. As an
effective consequence, if $T=a/b$ with $\log b\le800$, a single $h=3\ell-1\le36000$ with $\ell\nmid b$ gives a
contradiction. The exact scores $h^{-2}\log|F_h(T)|$ are much better than this bound, for example $-0.066$ at $h=192$
(Table~\ref{tab:evidence}).
\end{remark}

% ===========================================================================
\section{Two points}\label{sec:M1}
% ===========================================================================

\subsection{The statement}
Let $n\ge6$ be even and $P_n(X):=X^4+nX^3+2X^2-nX+1$. It is irreducible with four real roots: two in $(0,1)$, which we
call $z>w$, and $-1/z$, $-1/w$; moreover $w=(1-z)/(1+z)$, and $z\approx1-2/n$, $w\approx1/n$. Let $\cK:=\QQ(z)$, a totally
real biquadratic field with $\operatorname{Gal}(\cK/\QQ)=\{1,s,l,sl\}$, $s(z)=-1/z$, $l(z)=w$. Put $u:=z-1/z$; then
$u^2+nu+4=0$, and $\cF:=\cK^s=\QQ(u)=\QQ(\sqrt{n^2-16})$, on which $l$ acts as the nontrivial automorphism, with
$l(u)=4/u$. With $\eps_k:=(k+\sqrt{k^2-4})/2$ we have $zw=\eps_{n+2}^{-1}$ and $z/w=\eps_{n-2}$. Put $T_z:=\chi_2(z)$,
$T_w:=\chi_2(w)$; Landen's identity gives
\[
  T_n:=T_z+T_w=\frac{\pi^2}8-\frac12\log z\log w=\frac{\pi^2}8-\frac18\log^2\eps_{n+2}+\frac18\log^2\eps_{n-2}.
\]

\begin{definition}
$\mathrm{M1}(n)$ is the statement: $T_n\notin\QQ$, or $T_z-T_w\notin\sqrt{n^2-16}\,\QQ$.
\end{definition}

Equivalently, there is no $\alpha\in\cF$ with $T_z=\alpha$ and $T_w=l(\alpha)$: such an $\alpha$ would have
$\alpha+l(\alpha)\in\QQ$ and $\alpha-l(\alpha)\in\sqrt{n^2-16}\,\QQ$, and conversely. (Exchanging $z$ and $w$ only changes
the sign of $T_z-T_w$.)

\begin{theorem}[Theorem M1]\label{thm:M1}
$\mathrm{M1}(n)$ holds for every $n\in\mathcal N:=\{730,740,750,760,800,1000\}$.
\end{theorem}

The proof is computer-assisted: for each $n\in\mathcal N$ separately it uses two dual certificates checked in interval
arithmetic. It is a disjunction at six specific values of $n$, not a statement about a range of $n$, and it is not
formalised.

\emph{The case $n=730$} (Theorem~\ref{thm:M1-730}). Here $z=0.99726400650629232574\ldots$,
$w=0.0013698707255499998352\ldots$, $\eps_{732}=366+\sqrt{133955}=1/(zw)$, $\eps_{728}=364+11\sqrt{1095}=z/w$ and
$\sqrt{730^2-16}=22\sqrt{1101}$. Numerically $T_{730}=1.2246689334990568135\ldots$, $T_z=1.2232990624878816026\ldots$ and
$T_w=0.0013698710111752109\ldots$. The fields are $\cK=\QQ(\sqrt{133955},\sqrt{1095})$ and $\cF=\QQ(\sqrt{1101})$, and the
primes ramified in $\cK$ are $2,3,5,73,367$. The other instances are in Table~\ref{tab:M1fields}; in each case
$\cK=\QQ(\sqrt{d_+},\sqrt{d_-})$ with $d_\pm$ the squarefree parts of $(n\pm2)^2-4$.

\begin{table}[ht]
\centering\small
\begin{tabular}{@{}rllll@{}}
\toprule
$n$ & $\eps_{n+2}=1/(zw)$ & $\eps_{n-2}=z/w$ & $\sqrt{n^2-16}$ & ramified in $\cK$\\
\midrule
730 & $366+\sqrt{133955}$ & $364+11\sqrt{1095}$ & $22\sqrt{1101}$ & $2,3,5,73,367$\\
740 & $371+2\sqrt{34410}$ & $369+4\sqrt{8510}$ & $16\sqrt{2139}$ & $2,3,5,23,31,37$\\
750 & $376+5\sqrt{5655}$ & $374+5\sqrt{5595}$ & $2\sqrt{140621}$ & $2,3,5,13,29,373$\\
760 & $381+2\sqrt{36290}$ & $379+6\sqrt{3990}$ & $12\sqrt{4011}$ & $2,3,5,7,19,191$\\
800 & $401+20\sqrt{402}$ & $399+20\sqrt{398}$ & $4\sqrt{39999}$ & $2,3,67,199$\\
1000 & $501+10\sqrt{2510}$ & $499+10\sqrt{2490}$ & $4\sqrt{62499}$ & $2,3,5,83,251$\\
\bottomrule
\end{tabular}
\caption{The six instances of Theorem~\ref{thm:M1}.}\label{tab:M1fields}
\end{table}

\subsection{What changes at two points}
The machine is that of Section~\ref{sec:construction} with $z$ the root above: the same poles, atoms and functional,
now with $A,B\colon\QQ[t]\to\cK$ and $\Delta_h\in\cK[X]$. We write $y=z^2$; then $1-y=-uz$. Since $u$ is an algebraic integer
with $u\,l(u)=4$, $u$ divides $4$, and $z$, $w$ are units of $\cK$ (roots of a monic polynomial with constant term $1$).

\emph{Reflection.} Lemma~\ref{lem:R} holds with $\sigma=s$, the generator of $\operatorname{Gal}(\cK/\cF)$: its proof
used only $\sigma(z)=-1/z$ and that the coefficients are rational functions of $z$ with rational coefficients. So the
odd-index entries lie in $e\,\cF[X]$ with $e:=z+1/z$, $s(e)=-e$, $e^2=u^2+4=-nu\in\cF$, and the rescaled matrix
$E_{i+j}e^{\ind{i\text{ odd}}-\ind{j\text{ odd}}}$ shows $\Delta_h\in\cF[X]$. The reflection removes the conjugates
$-1/z$ and $-1/w$, where the series diverge.

\emph{The second place.} The remaining conjugate, $l$, is not removed but convergent: the coefficients of $A$ and $B$
are rational functions of $z$, so $l(\Delta_h)(T_w)=\det(\Lambda_w(t^{i+j}))$, the Hankel determinant of the same
machine at the weight $w$ (the series converge since $0<w<1$). The two embeddings of $\cF$ thus see the $z$-machine and
the $w$-machine, and the endgame uses the norm $\cF\to\QQ$.

\emph{Odd primes.} At every prime of $\cK$ above an odd $p$, $z$, $y$ and $1-y=-uz$ are units. These are the only
properties of $z$ that Lemmas~\ref{lem:atom}--\ref{lem:nine} and Theorem~\ref{thm:A} use, so they hold verbatim, with the
valuation normalised by $v(p)=1$; there are no special odd primes, including the ramified ones and those dividing
$n^2-16$. The bound is the same at every prime $\fp$ of $\cF$ above $p$.

\emph{The prime $2$.} Here Lemma~\ref{lem:ten} fails: at $n=400$, $h=48$, the exact $2$-adic content is $388.5$, above
$h\lfloor\log_2h\rfloor+3h=384$. It is replaced by the following weaker bound.

\begin{lemma}[Lemma~10$''$]\label{lem:tenpp}
Let $n$ be even and $\fP$ a prime of $\cK$ above $2$, with $v=v_\fP$ normalised by $v(2)=1$. Then
$-v(\Delta_h)\le(h-2)^2/4+3h$.
\end{lemma}

\begin{proof}
Since $n$ is even, $u/2$ is a root of $X^2+(n/2)X+1$, hence a unit; so
$\nu:=v(1-y)=v(u)=1$, instead of $2$. Follow the proof of Lemma~\ref{lem:ten}, which up to the valuations does not
involve $z$. In (L), $y^b/(1-y)^{b+1}$ now has valuation $-(b+1)$, so
$v(G_{ij})\ge1-\pi+\min_{b<\pi}(v_2(b!)-1)=-\pi$. In (U), with $v(\eta)=1$, the first part has valuation $2+s_2(N_0)$ and
the second part valuation $\ge1+s_2(N_0)$, so $v(G_{ij})\ge2$ for $N_0\ge0$; the bound no longer grows with $N_0$. Hence
$v(G_{ij})\ge-f(i+j)$ with $f(s):=(s-K)_+$, a convex function, and $v(\tilde\Delta_h)\ge-\max_\sigma\sum_if(i+\sigma(i))$.
If $i<i'$ and $\sigma(i)>\sigma(i')$, put $c=i+\sigma(i)$, $d=i'+\sigma(i')$, $a=i+\sigma(i')$, $b=i'+\sigma(i)$; then
$a+b=c+d$ and $a<c,d<b$, so convexity gives $f(c)+f(d)\le f(a)+f(b)$, and removing inversions does not decrease the sum.
So the identity maximises, and $\sum_{i<h}f(2i)=1+3+\dots+(h-3)=(h-2)^2/4$. Finally $-v(\Delta_h)=3h-v(\tilde\Delta_h)$.
\end{proof}

So the prime $2$ costs $\frac14\log2=0.1733\ldots$ per $h^2$ instead of $0$. The mechanism that made it free in
Lemma~\ref{lem:ten} is not available anywhere in the family:

\begin{proposition}\label{prop:nu}
For every $n$ there is a prime $\fP$ of $\cK$ above $2$ with $v_\fP(1-y)\le1$ (normalised by $v_\fP(2)=1$). In particular
$v_2(1-y)=2$, which holds at $\sqrt5-2$ and makes the prime $2$ free there, cannot hold at every prime above $2$.
\end{proposition}

\begin{proof}
$z$ is a unit, so $v_\fP(1-y)=v_\fP(u)$. From $u\,l(u)=4$ we get $v_\fP(u)+v_\fP(l(u))=2$, and
$v_\fP(l(u))=v_{l^{-1}(\fP)}(u)$. If $v_\fP(u)\ge2$ for every $\fP\mid2$, then $v_\fP(l(u))\le0$ for every $\fP$, so
$v_{l^{-1}(\fP)}(u)\le0$, a contradiction.
\end{proof}

For even $n$, $\nu=1$ at every prime above $2$; for odd $n$, $\nu\in\{0,2\}$, and the prime with $\nu=0$ is worse, so odd
$n$ are not used. Proposition~\ref{prop:nu} does not show that the loss $\frac14\log2$ is necessary: exact profiles suggest
that the true $2$-adic cost is about $\frac18\log2$ per $h^2$, but we have not proved this.

\emph{The analytic part.} Theorem~\ref{thm:B} holds at a general weight $q\in\{z,w\}$, with $L:=\log(1/q)$, the field
$v_q(x):=-Lx-\Psi(x)$, and three changes. (a) $v_q$ is strictly decreasing on $(-\tanh L,\infty)$, which contains
$[-1/h,\infty)$ if and only if $h>\coth L$; this is the range of $h$ ($\coth L_z=364.9986\ldots$ at $n=730$, and
$\coth L_w<1.0001$). (b) The dual density is only assumed $\ge0$, of mass $1$, with $\operatorname{supp}\mu\subset[0,X_e]$
for some $X_e\ge4$ (at the $z$-place the support is long, $X_e\approx1300$), and $\Phi_q(\mu,\lambda):=-E(\mu)+\lambda+
\int_0^{X_e}(F_\mu-\lambda)_+\,dx$ with $F_\mu=2U^\mu+v_q$. (c) The tail condition becomes: with
$B(x):=2\log x-\log(x-1)-Lx$, both $\beta:=\lambda-B(X_e)$ and $\gamma:=L-(X_e-2)/(X_e(X_e-1))$ are positive. With
$M_0:=2\log(X_e+1)+1+L-\lambda$ and $\kappa:=2/(1-e^{-2\gamma})$, the proof of Theorem~\ref{thm:B} gives, for every even
$h>\coth L$,
\[
  \log\bigl|\det(\Lambda_q(t^{i+j}))_{i,j<h}\bigr|\le h^2\Phi_q(\mu,\lambda)+E_q(h),
\]
\begin{multline*}
  E_q(h):=2h\log(4h)+h\log4+(hX_e/2+2)\log3\\
  +h\,c_1(h)+2.362h-h\log h+h\max(M_0,0)+\kappa .
\end{multline*}
The pattern count below $X_eh$ replaces $(2h+1)\log3$, the cell $[-1/h,0]$ costs $h\max(M_0,0)$ (here $U^\mu\le\log(X_e+1)$
and $v_q\le L+1$ there), and the far nodes cost $\kappa$. Since $E_q(h)=O(h\log h)$, the $\limsup$ is again $\Phi_q$.

\emph{Nonvanishing.} Lemma~\ref{lem:P} holds over $\cF$ for $h=3\ell-1$ with $\ell\ge7$ prime and unramified in $\cK$.
Now $\zeta:=z^\ell\bmod\fl$ is $\operatorname{Frob}_\fl(z)$, one of four values:
\begin{center}
\begin{tabular}{@{}ccccc@{}}
\toprule
$\operatorname{Frob}_\fl$ & $1$ & $s$ & $l$ & $sl$\\
\midrule
$\zeta$ & $z$ & $-1/z$ & $w$ & $-1/w$\\
$1-\zeta^2$ & $-uz$ & $u/z$ & $-(4/u)w$ & $(4/u)/w$\\
$1-\zeta^{-2}$ & $u/z$ & $-uz$ & $(4/u)/w$ & $-(4/u)w$\\
\bottomrule
\end{tabular}
\end{center}
Every entry is an $\fl$-unit because $u\mid4$ and $\ell$ is odd, and steps 1--5 of Lemma~\ref{lem:P} use only these units
and $1-z^{-2}=u/z$. Since $\ell$ is unramified, $v_\fl$ restricts to $v_\fL$ on $\cF$ for $\fL=\fl\cap\cF$, and step~6
gives $\Delta_h(\alpha)\ne0$ for every $\alpha\in\cF$ with $v_\fL(\alpha)\ge-1$, whether $\ell$ is split or inert in $\cF$.

\begin{lemma}[Lemma~E$'$, the norm endgame]\label{lem:Eprime}
Let $n\ge6$ be even, and suppose that for each $q\in\{z,w\}$ there is a pair $(\mu_q,\lambda_q)$ as in~(b) above
($\mu_q\ge0$, mass $1$, $\operatorname{supp}\mu_q\subset[0,X_e]$ with $X_e\ge4$; the cap $\mu_q\le1$ is not required)
satisfying the tail condition~(c), with
\[
  m(n):=\tfrac12\bigl(\Phi_z(\mu_z,\lambda_z)+\Phi_w(\mu_w,\lambda_w)\bigr)+\tfrac43+\tfrac14\log2<0 .
\]
Then $\mathrm{M1}(n)$ holds.
\end{lemma}

\begin{proof}
Suppose $\alpha\in\cF$ has $T_z=\alpha$ and $T_w=l(\alpha)$, and fix $b\ge1$ with $b\alpha\in\mathcal O_\cF$. Let
$\mathcal H:=\{3\ell-1:\ell\ge7\text{ prime, unramified in }\cK,\ \ell\nmid b,\ 3\ell-1>\coth L_z\}$, an infinite set; note
$\coth L_w<\coth L_z$. For $h\in\mathcal H$, Lemma~P$'$ over $\cF$ gives $\Delta_h(\alpha)\ne0$. For a prime $\fp$ of $\cF$
put $n_\fp:=-v_\fp(\Delta_h)$, which may be a half-integer at ramified primes, and $n_p:=\max(0,\max_{\fp\mid p}\lceil
n_\fp\rceil)$, $c_h:=\prod_pp^{n_p}$. Then $c_h\Delta_h\in\mathcal O_\cF[X]$, and by the odd-prime lemmas and
Lemma~\ref{lem:tenpp},
\[
  \log c_h\le\sum_{3\le p<K}B_p\log p+\bigl((h-2)^2/4+3h\bigr)\log2+\theta(K)+\log2,
\]
where $B_p$ are the bounds used in Theorem~\ref{thm:A}; the term $\theta(K)=O(h)$ pays for all ceilings. So
$\limsup_{h\in\mathcal H}h^{-2}\log c_h\le\frac43+\frac14\log2$. The algebraic integer $\beta:=b^hc_h\Delta_h(\alpha)$ is
nonzero. Its two real embeddings are $b^hc_h\Delta_h(T_z)=b^hc_h\det(\Lambda_z(t^{i+j}))$ and
$b^hc_h\,l(\Delta_h)(T_w)=b^hc_h\det(\Lambda_w(t^{i+j}))$, so
$N_{\cF/\QQ}(\beta)=b^{2h}c_h^2\det(\Lambda_z(t^{i+j}))\det(\Lambda_w(t^{i+j}))$ is a nonzero integer. Taking logarithms,
dividing by $2h^2$, and applying the analytic bound at both places,
$0\le\frac43+\frac14\log2+\frac12(\Phi_z+\Phi_w)=m(n)<0$, a contradiction.
\end{proof}

\subsection{The certificates}
For each $n\in\mathcal N$ and $q\in\{z,w\}$ the certificate is a piecewise constant $\mu_q\ge0$ on a grid of $[0,X_e]$
with mass exactly $1$, a rational $\lambda_q$, a rigorous Arb upper bound for $\Phi_q(\mu_q,\lambda_q)$, and an Arb check
of the tail condition with $X_e\ge4$; the logarithm $L=\log(1/q)$ is an Arb ball computed from the quartic, so every bound
encloses the value at the true $L$. Two independent evaluators were used, and where both certified the same $(n,q)$ the
smaller bound is taken. Table~\ref{tab:M1cert} lists the results; $\frac43+\frac14\log2=1.50662012847331966\ldots$, and
every $m(n)$ is computed in Arb from the printed balls and lies entirely below $0$.

\begin{table}[ht]
\centering
\begin{tabular}{@{}rccr@{}}
\toprule
$n$ & $\Phi_z\le$ & $\Phi_w\le$ & $m(n)\le$\\
\midrule
730 & $0.97551590832772$ & $-3.98960878950$ & $-0.000426312$\\
740 & $0.98154355971108$ & $-3.99641267680$ & $-0.000814430$\\
750 & $0.98568546476333$ & $-4.00312517276$ & $-0.002099725$\\
760 & $0.98610919549783$ & $-4.00974867868$ & $-0.005199613$\\
800 & $1.00091072617$ & $-4.03494034059$ & $-0.010394678$\\
1000 & $1.05765433525$ & $-4.14651125316$ & $-0.037808330$\\
\bottomrule
\end{tabular}
\caption{Certified bounds for Theorem~\ref{thm:M1} (upper endpoints of Arb balls, rounded up). At the $z$-place
$X_e=1300$ for $n\le760$, and $X_e\approx1202$ and $1503$ for $n=800$, $1000$; at the $w$-place $X_e\approx4$.}\label{tab:M1cert}
\end{table}

The constant $\Phi_z$ is positive because $z\approx1-2/n$ gives the $z$-machine almost no decay; the $w$-place, with
$w\approx1/n$, pays for it. The side constants are comfortable: at $n=730$, for example, $\beta=0.5719$,
$\gamma=0.001971$, $M_0\le11.17$, $\kappa\le508.4$ at the $z$-place, and $\beta=16.41$, $\gamma=6.426$ at the $w$-place.
A third evaluator, written independently of both, confirms $m(n)<0$ for all six $n$ and gives
$m(730)\le-0.000452$. As a validation, the three codes applied at Theorem~\ref{thm:main}'s weight $\sqrt5-2$ give
$-1.35926599$, $-1.35939764$ and $-1.35939884$, consistent with Theorem~\ref{thm:cert} and the primal value
$-1.3594016$ of Remark~\ref{rem:equil}.

\subsection{Scope}
We state the limits of Theorem~\ref{thm:M1} explicitly.
\begin{itemize}
\item It is a disjunction. We do not know which of the two statements holds, and neither is known by other means:
  $T_n\notin\QQ$ is a statement of strong-four-exponentials type, and we know of no result on $T_z-T_w$.
\item For very large $n$ a general criterion may apply: our evaluation of \cite[Thm.~1]{DHK2022b} with the shift
  $\frac12$ at $w^2$, at the real place where $w$ is small, suggests that its condition holds once
  $n\gtrsim7.6\cdot10^{13}$, which would give $\chi_2(w)\notin\cK$ and hence $\mathrm{M1}(n)$ for those even $n$. This
  evaluation is a sketch that has not been checked in detail or by a referee. For $730\le n\le1000$ the same quantity lies
  between $-12.7$ and $-12.5$, so Theorem~\ref{thm:M1} is not implied.
\item It is proved for the six listed values of $n$ only, each with its own certificates. Even $n$ between them are not
  covered, and odd $n$ are not covered at all: Lemma~\ref{lem:tenpp} needs $n$ even, and for odd $n$ one prime above
  $2$ has $v(1-y)=0$.
\item For all even $n\ge730$ the statement is supported only numerically. A proof would need interval certificates in $n$
  (the dual bound is monotone in $L$, but this needs to be written) and an argument as $n\to\infty$; neither exists.
\item The method cannot reach $n=720$: there the primal value of the equilibrium problem gives $m(720)\ge+0.00098$
  (computed numerically, with a density admissible up to rounding of order $10^{-12}$), so no dual certificate can
  succeed. Numerically the threshold of the method lies in $(720,730)$.
\item The theorem is not formalised in Lean. Its proof was checked by two AI referees and has not been read by a human.
\end{itemize}

% ===========================================================================
\section{The formalisation}\label{sec:lean}
% ===========================================================================

\subsection{What is formalised}
The Lean project \cite{PiTwoRepo} proves Theorem~\ref{thm:main} and the identity \eqref{eq:closedform} in Lean~4
\cite{LeanFour} (version v4.35.0) with Mathlib \cite{Mathlib} (version v4.35.0 \cite{Mathlib4}). It has 62 modules
under \lean{Chi2U\slash{}}, with about $16{,}700$ lines, of which about $7{,}800$ are ported verbatim from the
formalisation of \cite{Chi2Paper} \cite{Chi2Repo}: the prime number theorem, the interval arithmetic, the generic
analytic files, the criterion and the definition of $\chi_2$. All declarations are in the namespace \lean{Chi2U}. The
project is available at \url{https://github.com/telaviv/pi2-log2-golden-irrationality} \cite{PiTwoRepo}, together with this
paper and the scripts of Section~\ref{sec:comp}. Theorem~\ref{thm:M1} is not formalised.

\subsection{The statement}
The definitions that the reader has to check are in \lean{Defs.lean}, \lean{Statement.lean} and
\lean{Golden\slash{}Corollary.lean}:
\begin{quote}\ttfamily\footnotesize
noncomputable def chi2Term (z : $\mathbb{R}$) (n : $\mathbb{N}$) : $\mathbb{R}$ :=\\
\hspace*{2em}z \^{} (2 * n + 1) / (2 * (n : $\mathbb{R}$) + 1) \^{} 2\\[2pt]
noncomputable def chi2 (z : $\mathbb{R}$) : $\mathbb{R}$ := $\sum$' n : $\mathbb{N}$, chi2Term z n\\[2pt]
def TheoremU : Prop := Irrational (chi2 (Real.sqrt 5 - 2))\\[2pt]
theorem chi2\_sqrt5\_sub\_two\_irrational : TheoremU\\
theorem chi2\_sqrt5\_sub\_two\_eq :\\
\hspace*{2em}chi2 (Real.sqrt 5 - 2) = Real.pi \^{} 2 / 24 - 3 / 4 * Real.log goldenRatio \^{} 2\\
theorem golden\_irrational' :\\
\hspace*{2em}Irrational (Real.pi \^{} 2 - 18 * Real.log goldenRatio \^{} 2)
\end{quote}
The first theorem is in \lean{Main.lean}, the second in \lean{Golden\slash{}ClosedForm.lean}, the third in
\lean{Golden\slash{}Corollary.lean}. The definition of $\chi_2$ is identical to that of \cite{Chi2Repo}, and
\lean{Defs.lean} proves summability for $|z|\le1$, so \lean{chi2 (Real.sqrt 5 - 2)} is the genuine sum of the series.
\lean{goldenRatio} is Mathlib's $(1+\sqrt5)/2$ and \lean{Irrational} is Mathlib's predicate $x\notin\QQ$. The identity
\eqref{eq:closedform} is proved from Landen's identity for $\chi_2$, Euler's reflection formula and Landen's identity for
$\Li_2$ (\lean{Golden\slash{}Landen.lean}, \lean{Golden\slash{}Dilog.lean}), so the golden-ratio form is not a separate
assumption.

The command \lean{lake env lean Axioms.lean} prints
\begin{quote}\ttfamily\small
'Chi2U.chi2\_sqrt5\_sub\_two\_irrational' depends on axioms:\\
\hspace*{2em}[propext, Classical.choice, Quot.sound]\\
'Chi2U.golden\_irrational'' depends on axioms:\\
\hspace*{2em}[propext, Classical.choice, Quot.sound]
\end{quote}
and the same list for \lean{chi2\_sqrt5\_sub\_two\_eq} and for each of the main ingredients listed in \lean{Axioms.lean}.
These are Lean's three standard axioms. The development contains no \lean{sorry}, no added axiom, no
\lean{native\_decide}, no \lean{implemented\_by} and no \lean{@[extern]}, so no step relies on compiled code.

\subsection{Correspondence with the paper}
Table~\ref{tab:lean} lists the main formal counterparts. Paths are relative to \lean{Chi2U\slash{}}.

\begin{table}[ht]
\centering\small
\begin{tabular}{@{}>{\raggedright\arraybackslash}p{0.27\textwidth}>{\raggedright\arraybackslash}p{0.69\textwidth}@{}}
\toprule
paper & Lean declaration (file) \\
\midrule
$\KK$, $z$, $\sigma$ & \lean{Q5} $=$ \lean{QuadraticAlgebra} $\QQ$ $5$ $0$, \lean{zK}, \lean{sig}, \lean{emb} (\lean{Field})\\
Lemma~\ref{lem:pole}, \eqref{eq:mahler} & \lean{lemma1}, \lean{aeval\_chi2\_Delta} (\lean{PoleExpansion}); \lean{Lmom\_eq\_mahler} (\lean{Mahler})\\
Lemma~\ref{lem:R}, Cor.~\ref{cor:R} & \lean{eval$_2$\_Delta}, \lean{aeval\_Delta}, \lean{Delta\_map} (\lean{LemmaR})\\
Prop.~\ref{prop:criterion} & \lean{irrational\_of\_poly\_approx} (\lean{Criterion}); \lean{MachineBoundsU}, \lean{irrational\_of\_boundsU} (\lean{MainCond})\\
Lemma~\ref{lem:classcount} & \lean{AK\_eq} (\lean{Arith\slash{}ClassCount})\\
Lemma~\ref{lem:bordered} & \lean{lemmaB} (\lean{Arith\slash{}LemmaB})\\
Lemmas~\ref{lem:eight}, \ref{lem:nine} & \lean{lemma9} (\lean{Arith\slash{}OddPrimes}, \lean{Arith\slash{}Lemma9Levels}); \lean{Arith\slash{}Lemma8One}; \lean{padicNorm\_Delta\_odd} (\lean{Arith\slash{}OddExp})\\
Thm.~\ref{thm:A} & \lean{theoremA} (\lean{Arith\slash{}PrimeSums}); \lean{denAll\_int}, \lean{log\_denAll\_le} (\lean{Arith\slash{}Assembly})\\
Lemma~\ref{lem:ten} & \lean{lemma10\_L}, \lean{lemma10\_U}, \lean{padicNorm\_Delta\_two} (\lean{Arith\slash{}TwoAdic})\\
Thm.~\ref{thm:B} (other route) & \lean{log\_abs\_Dpoly\_ge}, \lean{DualCertU}, \lean{DualCertU.analytic} (\lean{Analytic\slash{}TheoremBU})\\
Thm.~\ref{thm:cert} (coarser) & \lean{exists\_certU} (\lean{Cert\slash{}CertU}); data \lean{Cert\slash{}DataU}\\
Lemma~\ref{lem:P} & \lean{lemmaP}, \lean{nonvan\_frequently} (\lean{NonVan\slash{}LemmaP})\\
Section~\ref{sec:proof} & \lean{boundsU}, \lean{chi2\_sqrt5\_sub\_two\_irrational} (\lean{Main})\\
\eqref{eq:closedform} & \lean{chi2\_landen}, \lean{Li2\_reflection}, \lean{Li2\_landen} (\lean{Golden\slash{}Landen}); \lean{chi2\_sqrt5\_sub\_two\_eq} (\lean{Golden\slash{}ClosedForm})\\
\bottomrule
\end{tabular}
\caption{Formal counterparts of the main steps.}\label{tab:lean}
\end{table}

\subsection{Differences from the written proof}
The formal proof establishes what the endgame needs, and in several places it takes a route that was shorter in Lean.
\begin{itemize}
\item \emph{The machine.} $\Delta_h$ is \emph{defined} as the determinant of the rational coordinates of the matrix
  $\tilde E$ of Corollary~\ref{cor:R}, so $\Delta_h\in\QQ[X]$ holds by definition; Lemma~\ref{lem:R} and
  Corollary~\ref{cor:R} then identify it with the Hankel determinant over $\KK$ and its value at $T$ with
  $\det(\Lambda(t^{i+j}))$.
\item \emph{Theorem B$'$.} The formal proof does not follow Section~\ref{sec:analytic}. It follows the lattice-energy
  route of the formal proof of Theorem~B in \cite{Chi2Paper}: the Andreief identity for the node expansion, the bound
  $h^h(1+2/h)^{h^2}$ for determinants with derivative rows via averages over roots of unity (instead of step~3), and a
  discrete P\'olya-type lattice energy inequality comparing the patterns, which live on the odd nodes, with a reference
  sequence on all of $\NN$ (instead of steps~5--6). The ingredient not present in \cite{Chi2Repo} is the two-sided Jensen
  bound of step~4
  (\lean{log\_abs\_Dpoly\_ge}), which gives the node weight \eqref{eq:nodeweight}. This route evaluates the dual at the
  nodes pointwise, so it never meets the cell $[-1/h,0]$ and does not use the monotonicity \eqref{eq:vmono} (which is
  nevertheless proved, as \lean{vU\_strictAntiOn}, on $[-\frac12,\infty)$). It is stated asymptotically, like Theorem~B of
  \cite{Chi2Paper}: the explicit $E_{\mathrm{an}}(h)$ is not formalised.
\item \emph{The certificate.} The formal certificate (\lean{DualCertU}) consists of a step density with $50$ cells of
  $[0,2]$ ($49$ nonzero), grid width $1/L$ with $L=200$ on $[0,X]$ with $X=4$, rational cell bounds $S_k$ for
  $2U^\mu+v$, and the tail conditions; its exponent is $-E(\mu)+\lambda+\frac1L\sum_k(S_k-\lambda)_+$. The cell bounds
  use chords and tangents of each term, as in \cite{Chi2Paper}, which keeps the cancellation between neighbouring cells
  without the near/far split of Section~\ref{sec:cert}. The certified value is
  $\Phi_\QQ=-54155169767979/40000000000000\approx-1.3538792$, coarser than Theorem~\ref{thm:cert} but enough: the
  assembly (\lean{boundsU}) uses $\beta=-27/20$, so the formal margin is $\frac43-\frac{27}{20}=-\frac1{60}$.
\item \emph{Lemma 8$'$.} Only the single-level range $K/3\le p$ used by Theorem~\ref{thm:A} is formalised, by a crude
  row bound on the Lagrange Gram matrix and Lemma~\ref{lem:bordered} instead of the atom expansion.
\item \emph{Lemma P$'$.} Congruences are taken in the residue field $\ZZ_{(\ell)}[\varphi]/\fl$ with $\zeta:=z^\ell$ itself;
  no identification of $\zeta$ by Frobenius is needed, only $1-\bar z^{2\ell}=(4\bar z)^\ell$. The progression sums are
  handled by an exact shift identity instead of the Newton form.
\item \emph{The endgame} uses only even $h$, through $P_h:=\Delta_{2\lfloor h/2\rfloor}$ in the criterion of
  \cite{Chi2Repo}.
\end{itemize}

\subsection{How the certificate is checked}
The certificate data (\lean{Cert\slash{}DataU.lean}) were produced by the script \lean{scripts\slash{}gen\_cert\_u.py}
(run as \lean{python scripts\slash{}gen\_cert\_u.py 50 8 4 Chi2U\slash{}Cert\slash{}DataU.lean}). The script is not trusted:
every check is a decidable statement about the rational data, proved by \lean{decide +kernel} in
\lean{Cert\slash{}CertU.lean} in $16$ chunks of grid cells, together with a table of $2004$ logarithm enclosures, the energy
bound and the tail. The dyadic enclosure of $\log2$ used by the table is itself checked by \lean{decide +kernel}. A wrong
value would make the build fail but could not make a false statement provable.

\subsection{Cost}
The kernel check of the certificate takes about $6.5$ minutes. A build of a fresh clone of the repository, with
Mathlib's prebuilt cache and at most four parallel Lean processes, took $9$ minutes $35$ seconds ($9110$ jobs, no errors)
on a Windows~11 desktop with $32$ logical processors and $31.7$\,GB of RAM, and \lean{Axioms.lean} then reported only the
three standard axioms.

\subsection{What has to be trusted}
Beyond the Lean kernel and the three axioms, the reader has to trust that the definitions displayed above express
Theorem~\ref{thm:main}, and the Mathlib definitions they use (real numbers, \lean{tsum}, \lean{Real.sqrt},
\lean{Real.log}, \lean{Real.pi}, \lean{goldenRatio}, \lean{Irrational}). Nothing else, in particular no numerical data and
no generating script, needs to be trusted. To avoid trusting the elaborator's handling of the environment, every
declaration of the $62$ project modules was replayed through the kernel with \lean{leanchecker} (shipped with the Lean
toolchain), each module against the compiled environment of its imports, on 7 November 2026; all modules passed, in
about $42$ minutes in total. Mathlib itself was not replayed from scratch (\lean{--fresh}). An AI referee checked the
fidelity of the statements (the definitions above, traced to Mathlib, and the absence of hypotheses, custom axioms and
compiler trust) and found them faithful. Apart from the check by the Lean kernel, all verification of the formalisation
so far was carried out by AI systems; no human expert has yet reviewed the Lean code.

% ===========================================================================
\section{Computations and reproducibility}\label{sec:comp}
% ===========================================================================

\subsection{Software}
The computations use python-flint~0.9.0 \cite{pythonflint}: exact arithmetic over $\QQ$ and $\QQ(\sqrt5)$ for every
statement about valuations, and Arb ball arithmetic \cite{Arb} for every real number that feeds a claim. Hankel
determinants are evaluated at $T$ in ball arithmetic with working precision about $3.6h^2+8000$ bits. The formalisation
uses Lean v4.35.0 and Mathlib v4.35.0.

\subsection{What each computation is used for}
\begin{itemize}
\item \emph{Proof inputs.} The certificate of Theorem~\ref{thm:cert} (\texttt{cert\_U.py}, run with arguments
  \texttt{1600 32 2 4}, $128$ bits, about one minute); the certificates of Theorem~\ref{thm:M1}
  (\texttt{dual\_ref.py}, and \texttt{dualcert\_arb2.py} with its library \texttt{dualcert.py}), the margins and the
  explicit numbers of Section~\ref{sec:M1}
  (\texttt{facts\_M1.py}, Arb at $3000$ bits). These are rigorous: floating point is used only to choose the density and
  $\lambda$, and any choice gives a valid bound.
\item \emph{The effective threshold} of Remark~\ref{rem:effective} (\texttt{h0\_U.py}, \texttt{h0\_scan.py}): exact sums,
  using the constants of Theorem~\ref{thm:B}.
\item \emph{Checks of the lemmas} (not used by the proof). The machine and the matrix $\tilde E$
  (\texttt{umach\_core.py}); the bordered identity at $h=2,\dots,12$ and $X\in\{0,1,-3/7\}$ (\texttt{bordered\_check.py});
  the row claim (\texttt{rowclaim\_check.py}); Lemmas~\ref{lem:eight}, \ref{lem:nine} on exact profiles up to $h=256$
  (\texttt{law\_check.py}, \texttt{u\_profile.py} with \texttt{fastcontent.py}); Lemma~\ref{lem:classcount} by brute
  force for every odd $d<K+40$ and $194$ values of $K$ (\texttt{law\_check.py}); Lemma~\ref{lem:ten} entrywise for $h\le64$ (\texttt{lemma10\_check.py}); Lemma~\ref{lem:P}
  (\texttt{blocks\_check.py}, \texttt{lemmaP\_check.py}). An AI referee reimplemented the machine without sharing code
  and reproduced the numbers it checked; the certificate reproduced bit for bit at $N=100,200,1600$.
\end{itemize}
These scripts are in the folder \texttt{scripts/paper} of the repository \cite{PiTwoRepo} (those for
Section~\ref{sec:M1} in \texttt{scripts/paper/m1}), with a README that gives the command line for each of them and the
small result files quoted here. Not included are the referees' own check scripts (apart from \texttt{dual\_ref.py}) and
the earlier script that produced the exact profiles for $h\le128$ in Table~\ref{tab:evidence}; those profiles are
included as data, and \texttt{u\_profile.py} computes the same quantities (checked at $h=16$).

\subsection{Numerical evidence}
Table~\ref{tab:evidence} shows exact scores. The arithmetic column is $h^{-2}\sum_pn_p\log p$ for the exact contents, so
the score is the sum of the two previous columns. The values approach the limits $-1.3592$ (or better) and
$\frac43$ slowly, as $O(\log h/h)$.

\begin{table}[ht]
\centering
\begin{tabular}{@{}rcccr@{}}
\toprule
$h$ & score $h^{-2}\log|F_h(T)|$ & analytic & arithmetic & $n_2$\\
\midrule
16 & $-0.2930$ & $-1.4379$ & $1.1449$ & 54\\
32 & $-0.1734$ & $-1.4195$ & $1.2462$ & 142\\
64 & $-0.1147$ & $-1.4015$ & $1.2869$ & 350\\
96 & $-0.0951$ & $-1.3930$ & $1.2980$ & 576\\
128 & $-0.0764$ & $-1.3876$ & $1.3112$ & 830\\
192 & $-0.0659$ & $-1.3815$ & $1.3156$ & 1344\\
256 & $-0.0596$ & $-1.3781$ & $1.3185$ & 1918\\
\bottomrule
\end{tabular}
\caption{Exact scores for Theorem~\ref{thm:main}. At $h=192$ and $256$ the contents are proved lower bounds for each
$n_p$ (fast content computation), which can only make the score look worse.}\label{tab:evidence}
\end{table}

\subsection{Controls}
No false instance of the method is available, since no unit with rational $\chi_2$ is known. As substitutes: the proved
finite-$h$ bound of Remark~\ref{rem:effective} lies above the exact score at every tested $h$, and every per-prime
inequality holds on the exact profiles, often with equality; for a unit of norm $+1$ the reflection breaks; with poles
on one side only the conjugate place costs $+0.50$ per $h^2$ at $h=16$; and at the points $\sqrt2-1$ and $\varphi^{-1}$ the
corresponding scores are positive (primal lower bounds $+0.3286$ and $+0.6572$ for the limiting exponent), so the
method does not conclude there. For Theorem~\ref{thm:M1}, the method does not conclude at $n=6$ ($m(6)\approx+0.477$),
$n=400$ ($m(400)\ge+0.070$) or $n=720$ (numerical values).

\section*{Acknowledgements and disclosure}
The research, the construction, the proofs, the certificate programs, the Lean formalisation, the literature searches and
the text of this paper were produced largely by AI models (Claude, Anthropic), working under the direction of the author. Apart from the check
by the Lean kernel, all verification so far, including the internal reviews of the argument by AI referees, was
carried out by AI systems. No human expert has yet reviewed the mathematics or the Lean code; we would be grateful for
corrections. The prime number theorem in the formalisation is ported from the PrimeNumberTheoremAnd project of
A.~Kontorovich, T.~Tao and collaborators \cite{PNTplus}, through the \lean{zeta5} repository of C\'esar Del Solar
\cite{zeta5} and the formalisation of \cite{Chi2Paper}.

\bibliographystyle{amsplain}
\bibliography{references}

\end{document}
